\chapter{动态与可变形 SLAM}
\label{ch:dynamic}

% ════════════════════════════════════════════════════════════════
%  新部「学习时代的 SLAM 与空间 AI」第二章。
%  吸收 SLAM Handbook ch15《Dynamic and Deformable SLAM》(Schmid, Montiel,
%  Huang, Cremers, Neira, Civera)：动态效应刻画(三轴框架)、动态物体剔除/跟踪
%  (动态物体因子图，对应 Handbook 式 15.1–15.4)、稠密动态、长期/change-aware
%  SLAM、可变形 SLAM(NRSfM/SfT、ED 图、DynamicFusion 非刚性融合)。
%  最高准则：本章是经典 SLAM「放松静态世界假设」后的自然延伸，读不出接缝——
%   * 开篇直接接住绪论里已埋的动态预告 sec:intro-dynamic(DynaSLAM/Co-Fusion/
%     DynamicFusion)，把「静态世界假设」何时失效讲清；
%   * 动态物体因子图写成本书 MAP 母方程 eq:intro-nls / 因子图 tab:factors 的
%     扩展(状态增广运动物体位姿 + 恒速因子)；
%   * 可变形稠密复用 ch:dense 的 TSDF(eq:dense-tsdf-fuse 的 W_eta 移动平均伏笔)
%     与 surfel(sec:dense-surfel 的非刚性形变图)，把刚性配准升级为非刚性；
%   * 形变能量泛函(ARAP/Killing)Handbook 全 punt → 本章自补其能量定义与正则项；
%     长/短期定义、local consistency 可识别性升级为本书自己的形式化论述。
%  右扰动为主线；Handbook 上下标式 T_i^w 一律改写为本书下标式 T_{wi}；
%  新符号(运动物体位姿/速度、形变场/节点权重)登记并避让全书已占符号。
% ════════════════════════════════════════════════════════════════

\begin{practice}[前置自测]
开始之前，先试答下列问题。它们勾勒了本章——也是经典 SLAM \emph{放松静态假设}后要回答的核心疑问。若已能流畅作答，可只速读本章的框架图与速查表；若卡住，正是本章要为你打通的。
\begin{enumerate}
  \item 经典 SLAM（\cref{ch:intro}）默认世界是\emph{静止刚体}。这条假设具体藏在哪两处？一旦场景里有走动的人、行驶的车，它会从哪一环开始把估计带偏？
  \item 同样是“运动的东西”，一个在相机眼前走过的人、和一把趁机器人离开房间被挪走的椅子，对 SLAM 是\emph{同一类}麻烦吗？区分它们的标准是绝对的时间长短，还是别的什么？
  \item 若不想丢弃动态物体、而要\emph{同时}估相机轨迹与每个运动物体的轨迹，需要往本书的因子图（\cref{tab:factors}）里新增哪些\emph{变量}和哪些\emph{因子}？“恒速运动模型”在其中扮演什么角色？
  \item 内窥镜下的器官、随风的旗帜——整个场景都在\emph{非刚性形变}、没有任何静止背景。此时连“相机相对什么定位”都成了问题。要让一团会变形的点云“站得住”，必须给它加什么样的\emph{先验}？
  \item 机器人第二次进同一个房间，发现角落里上次见过的椅子不见了。它怎么知道这是“椅子真被搬走了”，而不是“这次没往那个角落看”？这个看似哲学的区分，为什么会决定地图能不能被正确清理？
\end{enumerate}
\end{practice}

\section*{本章目标}
学完本章，你应当能够：
\begin{enumerate}
  \item \textbf{说清}经典 SLAM 的\emph{静态世界假设}藏在何处、何时失效，并用\emph{动态效应类型}/\emph{理解细节}/\emph{处理时机}三条轴刻画一个动态 SLAM 问题，据此区分\emph{状态估计}（把动态当外点）与\emph{场景表示}（重建运动与形变）两类目标；
  \item \textbf{形式化}短期动态与长期动态的区别——它由\emph{场景变化率与机器人观测率之比}决定，而非绝对时长——并据此把“剔除/跟踪/长期变化/可变形”四条线放进同一坐标；
  \item \textbf{推导}动态物体 SLAM 的因子图：把每个运动刚体的位姿与速度\emph{增广}进状态、把恒速运动写成因子，得到改进重投影残差（\cref{eq:dyn-reproj}）、恒速残差（\cref{eq:dyn-vcte}）与速度-位姿耦合残差（\cref{eq:dyn-couple}），并论证它仍是本书 MAP 母方程（\cref{eq:intro-nls}）的一个实例；
  \item \textbf{讲清}动态外点判别的统一原理（冗余 $+$ 鲁棒核 $+$ 多视一致性 $+$ 语义先验）与运动分割，理解“可能移动”不等于“正在移动”这一陷阱；
  \item \textbf{自补并写出}非刚性重建的\emph{形变能量泛函}——嵌入式形变图（ED graph）的蒙皮插值、as-rigid-as-possible（ARAP）正则、近等距/Killing 正则——把 \cref{ch:dense} 的刚性配准升级为非刚性配准（DynamicFusion 式），并说清单目可变形 SLAM 的\emph{浮动地图歧义}与其消解所需的先验；
  \item \textbf{论述}长期/change-aware SLAM 的核心可识别性原理——\emph{局部一致性}（local consistency），辨析\emph{“无证据” vs “有缺席证据”}、边缘化的“冻结”陷阱，并说明物体级表示与 $\SE(3)$-等变描述子为何在此举足轻重；
  \item \textbf{定位}本章在全书的位置：它是\cref{ch:learning} 之后“放松世界模型假设”的一章，向\cref{ch:semantic} 的语义、\cref{ch:openworld} 的开放世界、以及经典 \cref{ch:dense,ch:loop,ch:nlopt} 都留有接口。
\end{enumerate}

\subsection*{本章知识导航}
本章是一条“\emph{把经典 SLAM 的静态世界假设逐步放松，直到面对运动、形变、长期变化的真实世界}”的主线。前几部建立的因子图、TSDF/surfel、回环，都将在这里被\emph{推广}而非推翻。结构如下：

\begin{center}
\begin{forest} navtreeLR
[{放松静态假设\\三轴刻画}
  [{短期动态\\剔除 vs 跟踪}
    [{动态物体因子图\\(增广 + 恒速)}]]
  [{可变形 SLAM\\NRSfM / SfT}
    [{形变能量\\ED·ARAP·Killing}]
    [{非刚性融合\\(复用 ch:dense)}]]
  [{长期 / change-aware\\local consistency}
    [{地图清理\\变化检测}]
    [{时序场景理解\\(展望)}]]
]
\end{forest}
\end{center}

\noindent\textbf{推荐路径}：\cref{sec:dynamic-frame}（放松静态假设、三轴刻画）是全章的“地图”，把后面四条线放进同一坐标，任何读者都应先读；\cref{sec:dynamic-object}（短期动态与动态物体 SLAM）是本章\emph{唯一成体系的硬核数学}，含动态物体因子图的完整推导（接母方程），是重点自包含部分；\cref{sec:dynamic-deformable}（可变形 SLAM）体量最大、最全新，含本书\emph{自补}的形变能量泛函（ED/ARAP/Killing）与对 \cref{ch:dense} 稠密表示的非刚性升级；\cref{sec:dynamic-longterm}（长期/change-aware SLAM）给出一套关于“地图随时间变化”的概念框架与可识别性论证。带“（选读）”标记的稠密运动分割、医学内窥镜流水线、时序场景理解面向研究，可按需深潜。

\subsection*{前置知识桥接}
本章把前几部的工具“放松静态假设”后重新用一遍，请先激活以下前置：
\begin{itemize}
  \item \textbf{动态/可变形预告}（\cref{sec:intro-dynamic}）：绪论的前沿节已\emph{概览}预告了 DynaSLAM（剔除）、Co-Fusion（分别跟踪）、DynamicFusion（非刚性融合），并点破“动态 SLAM 把鸡生蛋循环又拧紧一圈”。本章正是\emph{兑现}那则预告——把概览升级为推导。
  \item \textbf{MAP 母方程与因子图}（\cref{eq:intro-nls};\cref{sec:est-fg},\cref{tab:factors}）：动态物体 SLAM \emph{不}另起炉灶——它只是把\emph{运动物体的位姿与速度}也当成因子图的变量、把\emph{运动模型}写成额外的因子，仍在解同一条母方程。
  \item \textbf{右扰动与 $\SE(3)$ 运动学}（\cref{ch:lie};\cref{ch:rigid_body}）：恒速运动模型 $\Delta\mathbf{T}=\exp((\boldsymbol\xi^k)^\wedge\Delta t)$ 与残差对位姿的求导，统一用本书右扰动 $\mathbf{T}\,\mathrm{Exp}(\delta\boldsymbol\xi)$（\cref{sec:est-fg-linear} 的线性化套路照搬）。
  \item \textbf{稠密表示：TSDF 与 surfel}（\cref{ch:dense};\cref{sec:dense-tsdf},\cref{sec:dense-surfel}）：可变形稠密重建\emph{直接复用}这里的体素融合与面元，只把\emph{刚性}配准换成\emph{非刚性}的形变图。\cref{ch:dense} 的 TSDF 加权移动平均（\cref{eq:dense-tsdf-fuse} 的 $W_\eta$ 封顶）已为“支持动态场景”埋下伏笔，ElasticFusion 的“无位姿图、直接形变面元地图”（\cref{sec:dense-surfel}）则是非刚性融合的雏形。
  \item \textbf{回环与地点识别}（\cref{ch:loop}）：长期定位的外观不变特征、物体实例重定位，与回环/地点识别\emph{同源}——都是“在变化中认出同一处地方”，本章交叉引用而非重述。
  \item \textbf{学习接进几何}（\cref{ch:learning}）：本章多处用到“网络出语义掩膜 / 网络出可变形先验 / 网络出 $\SE(3)$-等变描述子”，其方法论入口是 \cref{ch:learning} 已立的统一立场（把母方程的零件交给网络）。
\end{itemize}

\subsection*{如果跳过本章会怎样}
\begin{itemize}
  \item 你会默认 \cref{ch:intro}--\cref{ch:dense} 的整套 SLAM 只在“静态、刚体”的实验室里成立，一遇真实街道、家居、人体就束手无策，却不知道经典框架\emph{放松一条假设}就能延伸到这些场景；
  \item 你会以为“动态 SLAM”是一套与因子图无关的新算法，从而看不出 DynaSLAM/VDO-SLAM/DynamicFusion 仍在解本书的母方程（\cref{eq:intro-nls}），也学不会把运动物体或形变场\emph{接进}自己的优化问题；
  \item 在 \cref{ch:semantic}（语义为动态物体提供“它会不会动”的先验）与 \cref{ch:openworld}（开放世界里物体可被任意增删）中，你会反复遇到“场景在变”的设定——本章是它们共同的几何与概率基础。
\end{itemize}

\begin{note}
\textbf{本章记号与约定.}\;
本章在前几部记号上新增一层“动态/形变层”符号，并\emph{严格避让}全书已占用者，登记如下。
\textbf{相机位姿}沿用 $\mathbf{T}_{wi}\in\SE(3)$，表示把 $i$ 时刻相机系坐标变到世界系 $w$ 的位姿（\textbf{Handbook 原文的上下标式 $\mathbf{T}_i^w$ 一律改写为本书下标式 $\mathbf{T}_{wi}$}，见 \nameref{ch:notation}）；扰动\textbf{以右扰动为主线} $\mathbf{T}\,\mathrm{Exp}(\delta\boldsymbol\xi)$。
\textbf{运动物体}用上标 $k$ 标号（$k=1,\dots,K$ 为场景中的运动刚体）：第 $k$ 个物体在 $i$ 时刻的位姿记 $\mathbf{T}_{wi}^{k}\in\SE(3)$（把物体系 $k$ 在 $i$ 时刻的坐标变到世界系），其\emph{线速度}与\emph{角速度}记 $\mathbf{v}_i^{k},\boldsymbol\omega_i^{k}\in\mathbb{R}^3$（在物体系表达）；物体上的点在\emph{物体系}的坐标记 $\mathbf{p}_j^{k}\in\mathbb{R}^3$（不随时间变，刚体性），其齐次形式 $\tilde{\mathbf{p}}_j^{k}\in\mathbb{R}^4$。静态地图点（背景）沿用 $\mathbf{y}_l\in\mathbb{R}^3$（与 \cref{ch:intro} 一致，背景路标）。
\textbf{物体系恒速位姿增量}记 $\Delta\mathbf{T}^{k}_{i\to i+1}$。投影/反投影沿用 $\boldsymbol\pi,\boldsymbol\pi^{-1}$（\cref{ch:camera}）；齐次到欧氏的“去齐次”算子记 $h^{-1}:(x,y,z,\lambda)\mapsto(x/\lambda,y/\lambda,z/\lambda)$（仅本章局部用，勿与观测函数 $h$ 混；本章观测函数一律写全 $\boldsymbol\pi$）。
\textbf{可变形层}：标准模型（canonical map，即“静止时的形状”）记 $\mathcal{M}_c$；\textbf{形变场}（warp field）记 $\mathcal{W}:\mathbb{R}^3\to\SE(3)$，把标准模型的每个点送到当前帧；嵌入式形变图（ED graph）的\emph{节点}记 $\mathbf{g}_m\in\mathbb{R}^3$（$m=1,\dots,M$ 为形变节点），每节点配一个局部仿射/刚体变换 $(\mathbf{R}_m,\mathbf{t}_m)$ 与影响半径 $r_m$，节点对点 $\mathbf{p}$ 的\emph{蒙皮权重}记 $w_m(\mathbf{p})$（注意：这是\emph{形变节点权重}，与 \cref{ch:learning} 的逐像素置信 $w_{ij}$、\cref{ch:dense} 的过程噪声 $\mathbf{w}$ 不同，故带下标 $m$ 与显式自变量 $w_m(\mathbf{p})$）；点 $i$ 的空间邻居集记 $\mathcal{N}(i)$。形变能量记 $E_{\mathrm{def}}$，其数据项/正则项分别记 $E_{\mathrm{data}},E_{\mathrm{reg}}$。
\textbf{物体形状基}（主动形状模型）的系数记 $\mathbf{c}$、基向量记 $\mathbf{b}_k$（仅 \cref{sec:dynamic-object} 局部用）。本章不引入新的优化器记号，求解一律回 \cref{ch:nlopt}（GN/LM、Schur）。
\end{note}

% ════════════════════════════════════════════════════════════════
\section{放松静态世界假设：动态 SLAM 的刻画}
\label{sec:dynamic-frame}

\paragraph{动机：那条没说出口的假设。}
回到 \cref{ch:intro} 我们为机器人立下的两条基本方程——运动方程与观测方程（\cref{eq:slam-model}）——和那张五模块框架图（\cref{fig:intro-pipeline}）。它们看上去对“世界长什么样”一言未发，可\emph{整套经典 SLAM 都暗暗站在一条假设上：环境是静止的刚体}。这条假设并非可有可无的细节，它\emph{深深织进}了两个我们早已习以为常的环节\cite{schmid2026dynamic}：
\begin{itemize}
  \item \cref{ch:intro} 说过，若一个特征点\emph{停在原地不动}，那么\emph{重新观测它}就能直接约束相机位姿——回环（\cref{ch:loop}）正建立在“同一个地方还是同一个样子”之上；
  \item \cref{ch:dense} 说过，若底层环境\emph{保持不变}，把若干已知位姿下的观测\emph{融合}起来就能高效重建稠密模型——TSDF 的加权平均（\cref{eq:dense-tsdf-fuse}）正假定每个体素背后是\emph{同一}块表面。
\end{itemize}
真实世界却常常\emph{不}满足这条假设。一辆自动驾驶车要在城市里安全行驶，地图里除了静态背景，还有别的车、骑行者、行人在\emph{独立运动}；一台家用机器人建图时，人不停地在它周围走动、挪动物体、重新布置房间；更极端的，一台手术机器人在人体内航行，\emph{整个场景都在持续形变、没有任何静止背景}。这些都让“同时估计机器人状态与环境表示”变成远比静态情形棘手的问题。绪论的前沿节（\cref{sec:intro-dynamic}）已为此埋下预告，点了 DynaSLAM、Co-Fusion、DynamicFusion 三个名字，并留下一句洞见——\emph{动态 SLAM 把“鸡生蛋”循环又拧紧了一圈}。本章就把那则预告\emph{兑现}成推导。

\begin{insight}[一句话：动态 SLAM 是“鸡生蛋”再加一道关]
经典 SLAM 已要\emph{同时}估位姿与静态地图（\cref{sec:intro-chicken-egg}）。动态场景下还要\emph{同时}判断\textbf{每个观测到底属于静态背景、还是某个独立运动的物体}（甚至属于一团正在形变的表面）。这本质上是一个更难的\emph{数据关联}问题（\cref{sec:intro-frontback}）——把“看到了什么”进一步细分为“看到的东西在不在动、跟谁一起动”。理解了这一层，下面所有方法（剔除、跟踪、非刚性融合、变化检测）就都是\emph{回答这同一个问题}的不同侧面。
\end{insight}

\paragraph{从一个假设破裂讲起：动态点会“注入”错误运动。}
为什么动态物体非处理不可？设相机静止，眼前一辆车匀速驶过。落在车身上的特征点在图像里整体平移，几何上\emph{与“相机朝反方向平移、场景静止”完全无法区分}——这正是 \cref{thm:intro-scale} 单目尺度歧义同源的“运动归属歧义”。若后端把这些点当成静态背景喂进重投影因子，它就会\emph{误以为相机在动}，把车的运动“注入”自己的位姿估计：轻则精度下降，重则在动态占比高、遮挡大时跟踪发散。所以放松静态假设的第一步，是\emph{把“这个观测属于谁”显式地建模进来}。

\subsection{三条轴：刻画一个动态 SLAM 问题}
\label{sec:dynamic-axes}

\paragraph{为什么要先分类。}
动态效应五花八门——一个孤立刚体的平移、一个会变形的人、只在两次造访间发生的长期变化、乃至毫无静止背景的器官形变——若不先理出骨架，方法就会读成一份杂乱清单。理想的通用机器人当然要\emph{实时重建任意动态场景}，但问题的复杂度随所考虑效应的增多\emph{指数级}上升；幸而对具体任务，往往只需处理其中一个子集。动态效应可沿一套\emph{三轴}框架来刻画\cite{schmid2026dynamic}，本章即以此为组织骨架（\cref{tab:dyn-axes}）。

\begin{table}[H]\centering\small
\caption{刻画动态 SLAM 问题的三条轴（据 \cite{schmid2026dynamic}，本书重组）}
\label{tab:dyn-axes}
\begin{tabular}{p{0.20\textwidth} p{0.40\textwidth} p{0.32\textwidth}}
\toprule
轴 & 问的是什么 & 两端 / 本章落点 \\
\midrule
\textbf{动态效应类型} & 机器人的观测里出现了哪种动态？ & 静态背景 — 短期动态（在视野内运动）— 长期动态（视野外变化）— 可变形（无静止背景）\\
\textbf{理解的细节} & 任务要求机器人理解到多细？ & 只估\emph{机器人状态}（动态当外点，\cref{sec:dynamic-object}）— 估\emph{场景表示}（重建运动/形变，\cref{sec:dynamic-deformable}）— 估\emph{时序演化}（4D 理解，\cref{sec:dynamic-longterm}）\\
\textbf{处理的时机} & 必须实时，还是可离线？ & 在线（多数机器人应用）— 离线（如定期巡检的变化检测、批量 NRSfM）\\
\bottomrule
\end{tabular}
\end{table}

\paragraph{第一轴的关键：短期 vs 长期，是“相对率”而非绝对时长。}
三轴里最需厘清的是\emph{第一轴}对动态的细分。新手常以为“短期/长期”按时间长短分——眼前走过算短、隔天变化算长。更本质的划分不按绝对时长、而按\emph{变化率与观测率之比}\cite{schmid2026dynamic}；本书将其形式化如下。

\begin{definition}[短期动态与长期动态（相对率定义）]\label{def:dyn-short-long}
设场景的\emph{变化率}为 $r_{\mathrm{change}}$（单位时间内场景状态的改变量级），机器人的\emph{观测率}为 $r_{\mathrm{obs}}$（单位时间内对该处的观测次数）。几乎一切物理过程在足够短的时间尺度上都是\emph{连续}的；机器人感知到的是“短期”还是“长期”动态，\textbf{不取决于现象的绝对时长，而取决于二者之比} $r_{\mathrm{change}}/r_{\mathrm{obs}}$：
\begin{itemize}
  \item \textbf{短期动态}：$r_{\mathrm{change}}/r_{\mathrm{obs}}$ 小——机器人\emph{在视野内持续观测}到运动，变化呈\emph{平滑、连续}特征（如眼前滚动的球、走过的人）。可\emph{跟踪}。
  \item \textbf{长期动态}：$r_{\mathrm{change}}/r_{\mathrm{obs}}$ 大——运动\emph{发生在视野之外}，机器人只看到两次观测\emph{之间}累积的结果，变化呈\emph{突兀、跳变}特征（如趁机器人离开被挪走的瓶子，造访间发生的结构改变）。要\emph{重新识别与匹配}。
\end{itemize}
极端地，无静止背景、整场都在形变的情形（多属短期动态）称\emph{可变形 SLAM}。
\end{definition}

\begin{insight}[同一段运动，对不同机器人可短可长]
\cref{def:dyn-short-long} 有个直接推论：\emph{短期还是长期是观测者的属性，不是环境或物体的内在属性}\cite{schmid2026dynamic}。同一个物体沿同一条轨迹移动，对\emph{一直盯着它看}的机器人 A 是短期动态（连续轨迹），对\emph{看了一眼就走开、回来才发现它变了位置}的机器人 B 却是长期动态（瞬移般的跳变）。在多数机器人场景里，人或物的运动速率（$\sim 10^{-1}$ 到 $10^{2}$ m/s）与传感器观测率（$\sim 10^{0}$ 到 $10^{3}$ Hz）量级相当，于是这条判据常\emph{简化}为：\textbf{现在是不是在视野里看着它动}——看着，短期；没看着、回来发现变了，长期。这也解释了为何多机器人/多次造访（\cref{sec:dynamic-longterm}）天然是长期动态的温床：造访\emph{之间}机器人不观测，变化在那时累积。
\end{insight}

\paragraph{第二轴：估状态，还是估表示。}
正如 \cref{ch:dense} 指出 SLAM 是“定位”与“建图”的二元过程，动态场景同样如此。若任务\emph{只需机器人状态}（如纯导航的车，只要知道自己在哪），把动态观测当\emph{外点或噪声}剔除、只对静态背景定位，往往就够了——这与静态 SLAM 用 RANSAC/鲁棒核抗外点（\cref{ch:nlopt}）一脉相承，我们在 \cref{sec:dynamic-object} 展开剔除与跟踪。若任务\emph{要稠密地图}（如要与物体交互的操作机器人），就得\emph{检测并表示}场景中的部分或全部运动与形变——这就走向稠密动态（\cref{sec:dynamic-object} 末）与可变形重建（\cref{sec:dynamic-deformable}）。一个特例是可变形 SLAM：\emph{没有静止部分可供定位}，状态估计与表示被迫纠缠在一起求解。

\paragraph{第三轴：在线还是离线。}
多数机器人要\emph{实时}与动态环境交互，故本章主讲在线方法。但某些应用（如定期巡检一栋楼）允许\emph{先采后算}：离线处理多一份算力、且\emph{全部未来数据已在手}——要判断 $t$ 时刻某物是否运动/变化，未来帧 $t+1,\dots,T$ 都可用，这是在线方法没有的信息。这一差别在长期变化检测（\cref{sec:dynamic-longterm}）与批量 NRSfM（\cref{sec:dynamic-deformable}）里尤为关键。

\begin{insight}[本章不按传感器、而按动态特征组织]
\cref{ch:camera}--\cref{ch:lidar_slam} 按\emph{传感器模态}（视觉、激光、雷达、IMU）组织。但动态场景对 SLAM 的影响\emph{在各模态间高度相似}（都要回答“谁在动”），故本章\emph{改按动态特征}（短期/长期/可变形）组织，每节再点出不同模态的考量——目前视觉、激光、本体感受传感器在动态/可变形 SLAM 用得最多。这与 \cref{ch:learning} 按“介入深度”组织、而非按系统罗列，是同一种“抓主线、避清单”的写法。
\end{insight}

% ════════════════════════════════════════════════════════════════
\section{短期动态：从剔除外点到跟踪运动物体}
\label{sec:dynamic-object}

本节处理\emph{短期动态}——机器人视野内独立运动的物体。历史上的第一反应是\emph{数据关联 $+$ 外点剔除}：用统计检验、RANSAC（\cref{ch:vo}）、霍夫变换等手段，把“能被相机自身运动解释”的观测与“呈独立运动”的观测分开，将后者当作刚性模型的\emph{虚假测量}丢弃\cite{schmid2026dynamic}。后来人们不再满足于丢弃，而是进一步\emph{修补}被动态物遮挡的背景（插值、从别的位姿补全、或用网络\emph{图像补绘} inpaint），得到一张无“空洞”、更利于定位的干净地图——这正是 \cref{sec:intro-dynamic} 预告的 DynaSLAM\cite{bescos2018dynaslam} 那一路。另一条路则\emph{正面跟踪}动态物体，把它纳入 SLAM 框架，在计算机视觉里称\emph{多体运动恢复结构}（multi-body SfM）——这也是本节的数学重头戏，因为它能与本书的因子图无缝衔接。

\subsection{动态外点剔除：把动态当虚假测量}
\label{sec:dynamic-removal}

\paragraph{统一原理：对刚性模型的“不一致”。}
在部分动态的场景里，最简单也最有效的一招，就是\emph{丢弃落在动态部分上的观测、并把它们从地图状态里移除}\cite{schmid2026dynamic}。理由有二：其一，持久的物体与路标往往才是建图的重点，路过的行人通常无足轻重；其二，几乎所有几何估计技术都\emph{假定场景刚性}，于是动态部分天然表现为对刚性模型的\emph{虚假测量}。判别“虚假”的手段众多，现代系统常组合使用，但背后是\emph{同一条原理}——\textbf{在“场景刚性”这个零假设下，动态观测是不一致的外点}。我们把主要手段归纳如下，不逐一复述各系统的流水线。

\begin{itemize}
  \item \textbf{本体感受去歧义}。惯性、轮速等本体感受测量\emph{不依赖外部观测}，故\emph{不受场景运动影响}——它们直接携带机器人自身的运动信息，可用来把“相机的运动”与“场景里物体的运动”从相对测量中\emph{解耦}\cite{schmid2026dynamic}。如何融合惯性/运动学与外感测量，见 \cref{ch:vio,ch:leg}。
  \item \textbf{鲁棒估计（冗余）}。多数传感器提供的测量\emph{远多于}几何估计所需的最小集。在动态占比不高时，可用\emph{冗余}把动态观测识别为外点：其一，用 RANSAC（\cref{ch:vo}）选\emph{最大一致集}初始化（但大遮挡下最大一致集可能反被运动物体占据，会失效）；其二，迭代优化时把 $\ell_2$ 范数换成\emph{鲁棒核}（\cref{ch:nlopt}）——它对大残差\emph{次二次}增长，从而压低疑似外点的影响。
  \item \textbf{多视运动线索}。除了在初始化时剔、用鲁棒核压，还可在\emph{累积到足够非刚性证据}后\emph{移除}动态状态——证据通常实现为“高残差测量累计数”的阈值。例如把 RGB-D 图聚成区域、按各区域的\emph{配准残差}把它二分为动态/静态背景，再用这一分割反过来\emph{精化}里程计（如 \cite{jaimez2017fast} 的联合相机运动 $+$ 场景流估计）。
  \item \textbf{语义先验}。用语义分割网络（\cref{ch:learning}）\emph{预先}标出“可能移动”的类别（人、车）。但要警惕：闭集词典未必能分出图中所有物体；更要紧的是——\textbf{“可能移动”不等于“正在移动”}（一辆停着的车是潜在动态、却是当下的好路标）。故语义掩膜常需与\emph{多视几何检验}联用\cite{bescos2018dynaslam}。MonoRec\cite{wimbauer2021monorec}（\cref{ch:learning} 已提）即把亮度一致性代价体与一个\emph{掩膜模块}结合——掩膜既从训练数据、也从\emph{扭曲帧间的颜色不一致}识别并滤除运动物。
  \item \textbf{领域知识}。可注入先验约束。如自动驾驶的“无侧滑约束”：车不能横向平移，故一切暗示侧向运动的特征可判为动态外点而剔除。
  \item \textbf{运动感知传感器}。普通视觉/激光只给某一时刻的快照、不直接测物体运动；而事件相机（\cref{ch:event}）给出高时间分辨率的连续测量流，雷达（\cref{ch:radar}）与某些激光用多普勒效应\emph{直接测每点的径向速度}——这类信息是识别动态物及其运动的有力线索。
\end{itemize}

\begin{pitfall}[陷阱：把“语义可动”直接当“动态外点”剔掉]
\textbf{现象}：用 Mask R-CNN 把所有“人、车”像素一律剔除，结果在停车场里把\emph{一排静止的车}全删了，背景路标骤减、定位反而变差。\textbf{根本原因}：语义只回答“这类东西\emph{会不会}动”，不回答“它\emph{此刻}动没动”。\textbf{正确做法}：语义掩膜只作\emph{候选}，再叠加多视几何一致性检验（落在掩膜内\emph{且}残差异常的才判动态）——这正是 DynaSLAM\cite{bescos2018dynaslam} “语义 $+$ 多视几何”双保险的由来。\textbf{自检}：你的动态判据里，有没有一条是\emph{基于当前帧的几何不一致}、而非仅凭类别？
\end{pitfall}

\subsection{动态物体跟踪：把运动刚体增广进因子图}
\label{sec:dynamic-tracking}

\paragraph{动机：有时必须知道“那辆车往哪开”。}
剔除够用于\emph{只要机器人状态}的场合，但很多任务\emph{必须}跟踪动态物体的三维运动——比如为了避免与它相撞。这一问题也被称作\emph{多目标跟踪}（MOT）或\emph{多实例动态 SLAM}（MID）。其文献浩如烟海，本节只取\emph{与 SLAM 最相关}的提法：把“同时定位、建图、跟踪动态物体”统一成\emph{一张因子图}\cite{schmid2026dynamic}。中心思想极朴素，也极契合本书主线：

\begin{insight}[核心思想：每个运动物体 $=$ 一条额外的刚体轨迹]
把场景里每个独立运动的物体都\emph{建模为一个刚体}。于是 \cref{sec:est-fg} 那张静态 SLAM 因子图就能\emph{自然扩展}：相机位姿之间用\emph{里程计因子}相连（不变），\textbf{每个动态物体的位姿序列之间用\emph{恒速运动因子}相连}，物体上的点经“先把物体摆到世界、再投到当前相机”产生改进的重投影因子。换言之——\textbf{动态物体 SLAM 没有跳出本书的母方程（\cref{eq:intro-nls}），它只是往状态里\emph{增广}了运动物体的位姿与速度，往因子图里\emph{添}了运动模型这一类因子}。这与 \cref{ch:learning} 把网络“接进”母方程、与 \cref{sec:intro-multirobot} 把多机器人轨迹塞进同一张大因子图，是\emph{同一种}“扩展而非重写”的手法。
\end{insight}

\paragraph{状态与变量。}
沿用本章记号（见章首 note）。设：
\begin{itemize}
  \item 相机位姿 $\mathbf{T}_{wi}\in\SE(3)$：$i$ 时刻相机系到世界系 $\mathcal{F}^w$ 的刚体变换；
  \item 静态背景点 $\mathbf{y}_l\in\mathbb{R}^3$：在世界系表达的刚性地图点 $l$；
  \item 动态物体位姿 $\mathbf{T}_{wi}^{k}\in\SE(3)$：把第 $k$ 个运动物体\emph{在 $i$ 时刻}的物体系 $\mathcal{F}^k$ 坐标变到世界系；
  \item 物体的线/角速度 $\mathbf{v}_i^{k},\boldsymbol\omega_i^{k}\in\mathbb{R}^3$：$i$ 时刻在物体系表达；
  \item 物体上的点 $\mathbf{p}_j^{k}\in\mathbb{R}^3$：在\emph{物体系}表达。\textbf{关键}：刚体运动学下，点在物体系的坐标\emph{不随时间变}，故 $\mathbf{p}_j^{k}$ 不带时间下标 $i$——这正是“刚体”二字的数学含义，也是把物体当刚体的全部好处所在。
\end{itemize}

\paragraph{因子一：改进的重投影残差。}
静态点的重投影是“把世界点投到当前相机”。动态点多一步：先用物体位姿把\emph{物体系}的点摆到\emph{世界系}，再用相机位姿的逆投到\emph{相机系}、最后投影。于是点 $\mathbf{p}_j^{k}$ 在 $i$ 时刻的预测像素与实测像素 $\mathbf{z}_j^{i}$ 之差给出
\begin{equation}
  \mathbf{e}_{\mathrm{reproj}}^{i,j,k}
  =\mathbf{z}_j^{i}-\boldsymbol\pi\!\Big(\,\mathbf{T}_{wi}^{-1}\,\mathbf{T}_{wi}^{k}\,\tilde{\mathbf{p}}_j^{k}\,\Big),
  \label{eq:dyn-reproj}
\end{equation}
其中 $\tilde{\mathbf{p}}_j^{k}\in\mathbb{R}^4$ 是 $\mathbf{p}_j^{k}$ 的齐次坐标、$\boldsymbol\pi(\cdot)$ 为针孔投影（\cref{ch:camera}，此处对齐次点重载）。

\begin{insight}[\cref{eq:dyn-reproj} 与静态重投影逐项对照]
把 \cref{eq:dyn-reproj} 与 \cref{tab:factors} 的静态观测残差 $\mathbf{z}_{k,j}-h(\mathbf{x}_k,\mathbf{y}_j)$ 并排看：静态情形里地图点 $\mathbf{y}_l$ 直接活在\emph{世界系}（不需要 $\mathbf{T}_{wi}^{k}$ 这一步）；动态情形里点活在\emph{物体系}，须先被 $\mathbf{T}_{wi}^{k}$ 送进世界。\textbf{两者唯一的差别，就是中间多插了一个“物体位姿” $\mathbf{T}_{wi}^{k}$}。令 $\mathbf{T}_{wi}^{k}\equiv\mathbf{I}$（物体不动、就是背景），\cref{eq:dyn-reproj} 立刻退回静态重投影。这就是“动态物体 SLAM 是静态 SLAM 的推广、而非另一套”的最直白证据。残差对 $\mathbf{T}_{wi}$ 与 $\mathbf{T}_{wi}^{k}$ 的雅可比，照 \cref{sec:est-fg-linear} 用右扰动 $\mathbf{T}\,\mathrm{Exp}(\delta\boldsymbol\xi)$ 链式求导即可，与 \cref{ch:vo} 重投影对相机位姿求导同构。
\end{insight}

\paragraph{因子二：恒速运动模型。}
只有重投影还不够：单帧里物体的位姿与点深度存在歧义（与单目尺度歧义同源），需要\emph{时间上的约束}把物体的运动“串”起来。最常用的是\emph{恒速模型}——假定物体在相邻观测间线/角速度近似不变。这给出一条直接的速度残差
\begin{equation}
  \mathbf{e}_{\mathrm{vcte}}^{i,k}
  =\begin{bmatrix}\mathbf{v}_{i+1}^{k}-\mathbf{v}_i^{k}\\[1mm]\boldsymbol\omega_{i+1}^{k}-\boldsymbol\omega_i^{k}\end{bmatrix}.
  \label{eq:dyn-vcte}
\end{equation}
它在精神上与 \cref{tab:factors} 的\emph{里程计因子}同型：里程计约束相邻\emph{相机}位姿之差，恒速因子约束相邻\emph{物体}速度之差。

\paragraph{因子三：速度-位姿耦合。}
速度残差 \cref{eq:dyn-vcte} 单独并不能把速度与位姿挂钩。还需一项\emph{耦合}残差，要求“由运动模型推得的物体点位置”与“由位姿估计推得的位置”一致。先由恒速模型写出物体从 $t_i$ 到 $t_{i+1}$ 的\emph{位姿增量}：在物体系里，角速度 $\boldsymbol\omega_i^{k}$ 与线速度 $\mathbf{v}_i^{k}$ 在 $\Delta t=t_{i+1}-t_i$ 内积分为
\begin{equation}
  \Delta\mathbf{T}^{k}_{i\to i+1}
  =\begin{bmatrix}\exp\!\big((\boldsymbol\omega_i^{k})^\wedge\,\Delta t\big) & \mathbf{v}_i^{k}\,\Delta t\\[1mm] \mathbf{0}^\top & 1\end{bmatrix},
  \label{eq:dyn-increment}
\end{equation}
其中 $(\cdot)^\wedge$ 是 $\so(3)$ 的反对称化（\cref{ch:lie}），$\exp$ 即 $\SO(3)$ 指数映射 $\mathrm{Exp}$。\cref{eq:dyn-increment} 是一个一阶（匀速）刚体运动的离散积分；本书\emph{统一}用 $\mathrm{Exp}/\mathrm{Log}$ 与右扰动记号，故把 Handbook 原文的 $\exp(\boldsymbol\omega(t_{i+1}-t_i))$ 写成上式更显式的 $\SE(3)$ 块形（旋转部分走 $\SO(3)$ 指数、平移部分按匀速 $\mathbf{v}\,\Delta t$，这是“匀速螺旋运动”的一阶近似，详见 \cref{ch:lie} 的指数映射）。耦合残差则比较“按位姿摆放”与“按运动模型摆放”同一物体点的世界坐标：
\begin{equation}
  \mathbf{e}_{\mathrm{couple}}^{i,j,k}
  =h^{-1}\!\big(\mathbf{T}_{w,i+1}^{k}\,\tilde{\mathbf{p}}_j^{k}\big)-h^{-1}\!\big(\mathbf{T}_{wi}^{k}\,\Delta\mathbf{T}^{k}_{i\to i+1}\,\tilde{\mathbf{p}}_j^{k}\big),
  \label{eq:dyn-couple}
\end{equation}
其中 $h^{-1}:(x,y,z,\lambda)\mapsto(x/\lambda,y/\lambda,z/\lambda)$ 把齐次坐标去齐次（见章首 note）——它\emph{分别}作用于两个变换后的点（各自齐次分量均为 $1$、良定义），再在欧氏空间作差；切勿先对两位姿之差作用 $h^{-1}$（其齐次分量为 $0$、未定义）。直觉上，$\mathbf{T}_{wi}^{k}\,\Delta\mathbf{T}^{k}_{i\to i+1}$ 是“从 $i$ 时刻位姿出发、按匀速走一步”应到达的 $i+1$ 时刻位姿；若它与独立估计的 $\mathbf{T}_{w,i+1}^{k}$ 作用在同一点上结果一致，残差就为零。

\begin{insight}[动态物体因子图 $=$ 母方程 $+$ 三类新因子]
把三类残差并进本书母方程，动态物体 SLAM 的 MAP 目标即
\begin{equation}
  \min_{\{\mathbf{T}_{wi}\},\{\mathbf{y}_l\},\{\mathbf{T}_{wi}^{k},\mathbf{v}_i^{k},\boldsymbol\omega_i^{k},\mathbf{p}_j^{k}\}}
  \sum_{\text{背景}}\!\big\|\mathbf{e}_{\mathrm{reproj}}^{\text{static}}\big\|^2_{\boldsymbol\Sigma}
  +\sum_{\text{里程计}}\!\big\|\mathbf{e}_{\mathrm{odom}}\big\|^2_{\boldsymbol\Sigma}
  +\underbrace{\sum_{i,j,k}\!\big\|\mathbf{e}_{\mathrm{reproj}}^{i,j,k}\big\|^2_{\boldsymbol\Sigma}
  +\sum_{i,k}\!\big\|\mathbf{e}_{\mathrm{vcte}}^{i,k}\big\|^2_{\boldsymbol\Sigma}
  +\sum_{i,j,k}\!\big\|\mathbf{e}_{\mathrm{couple}}^{i,j,k}\big\|^2_{\boldsymbol\Sigma}}_{\text{动态新增的三类因子}} .
  \label{eq:dyn-map}
\end{equation}
前两项原封不动是 \cref{eq:intro-nls} 的静态 SLAM（背景重投影 $+$ 里程计）；花括号里的三项就是动态物体带来的全部新增——\emph{每个运动物体一条轨迹、运动当因子}。这张图依旧\emph{稀疏}（每个物体只连自己的位姿/速度/点），故仍可用 \cref{ch:nlopt} 的 GN/LM $+$ Schur 消元高效求解。
\end{insight}

\begin{figure}[ht]
\centering
\begin{tikzpicture}[
  >=Stealth, node distance=12mm,
  cam/.style={circle, draw=main!70!black, fill=main!8, minimum size=8mm, font=\small, inner sep=0pt},
  obj/.style={circle, draw=third!70!black, fill=third!12, minimum size=8mm, font=\small, inner sep=0pt},
  vel/.style={circle, draw=third!55!black, fill=third!6, minimum size=7mm, font=\scriptsize, inner sep=0pt},
  fac/.style={rectangle, draw=second!60!black, fill=second!60!black, minimum size=2.1mm, inner sep=1.6pt},
  ed/.style={gray!55!black}]
  % camera poses
  \node[cam] (c1) {$\mathbf{T}_{w1}$};
  \node[cam, right=16mm of c1] (c2) {$\mathbf{T}_{w2}$};
  \node[cam, right=16mm of c2] (c3) {$\mathbf{T}_{w3}$};
  % odometry factors
  \node[fac] (od1) at ($(c1)!0.5!(c2)$) {}; \draw[ed] (c1)--(od1) (od1)--(c2);
  \node[fac] (od2) at ($(c2)!0.5!(c3)$) {}; \draw[ed] (c2)--(od2) (od2)--(c3);
  % static landmark
  \node[obj, above=11mm of c2, draw=main!60!black, fill=main!6] (y) {$\mathbf{y}_l$};
  \node[fac] (my1) at ($(c1)!0.5!(y)$) {}; \draw[ed] (c1)--(my1) (my1)--(y);
  \node[fac] (my2) at ($(c2)!0.5!(y)$) {}; \draw[ed] (c2)--(my2) (my2)--(y);
  % dynamic object poses (below)
  \node[obj, below=14mm of c1] (o1) {$\mathbf{T}_{w1}^{k}$};
  \node[obj, below=14mm of c2] (o2) {$\mathbf{T}_{w2}^{k}$};
  \node[obj, below=14mm of c3] (o3) {$\mathbf{T}_{w3}^{k}$};
  % reprojection factors camera<->object
  \node[fac] (r1) at ($(c1)!0.5!(o1)$) {}; \draw[ed] (c1)--(r1) (r1)--(o1);
  \node[fac] (r2) at ($(c2)!0.5!(o2)$) {}; \draw[ed] (c2)--(r2) (r2)--(o2);
  \node[fac] (r3) at ($(c3)!0.5!(o3)$) {}; \draw[ed] (c3)--(r3) (r3)--(o3);
  % constant-velocity (coupling) factors between object poses
  \node[fac] (v1) at ($(o1)!0.5!(o2)$) {}; \draw[ed] (o1)--(v1) (v1)--(o2);
  \node[fac] (v2) at ($(o2)!0.5!(o3)$) {}; \draw[ed] (o2)--(v2) (v2)--(o3);
\end{tikzpicture}
\caption{动态物体跟踪的 SLAM 因子图（仿 \cite{schmid2026dynamic,bescos2018dynaslam} 重绘）。顶行为相机位姿 $\mathbf{T}_{wi}$（蓝圆），以\emph{里程计因子}相连；上方为静态背景点 $\mathbf{y}_l$，以\emph{观测/重投影因子}连相机。底行为同一动态物体的位姿序列 $\mathbf{T}_{wi}^{k}$（橙圆），以\emph{恒速/耦合因子}（\cref{eq:dyn-vcte,eq:dyn-couple}）相连，并经\emph{改进重投影因子}（\cref{eq:dyn-reproj}）连到对应时刻的相机。测量不画为节点（\cref{def:factor-graph}）。与 \cref{fig:factor-graph} 的静态因子图相比，仅多了“底行物体轨迹 $+$ 恒速因子”一整条。}
\label{fig:dyn-factorgraph}
\end{figure}

\paragraph{代表系统与铰接体推广。}
这一通用提法已在多项工作中成功落地，本书只点其代表、不复述工程管线\cite{schmid2026dynamic}：早期的\emph{多运动视觉里程计}（MVO）检测轨迹片并聚成满足刚体约束的簇；\emph{VDO-SLAM}、Dynamic SLAM 与 \emph{DynaSLAM II}\cite{bescos2018dynaslam} 引入上述因子图提法、在观测间加运动一致性因子；\emph{AirDOS} 更进一步，把它推广到由\emph{多个相连刚体}组成的\textbf{铰接物体}（如人体骨架）——每段肢体一个刚体、关节处加约束，本质是“多个 \cref{eq:dyn-map} 式刚体 $+$ 关节因子”的组合。

\subsection{稠密动态：连运动物体也重建（选读）}
\label{sec:dynamic-dense}

\paragraph{从“跟踪”到“跟踪 $+$ 稠密重建”。}
若任务（如操作、人机交互）不仅要物体的\emph{位姿}、还要其\emph{稠密形状}，就走向\emph{稠密动态 SLAM}\cite{schmid2026dynamic}。多数方法把过程拆成\emph{定位}（对静态背景跟踪相机，复用本节剔除/MOT 的任意一招）与\emph{建图}（为背景\emph{与}每个动态实体各建一个稠密模型）两步。背景重建与静态稠密（\cref{ch:dense}）完全相同；动态实体则常用\emph{同样}的表示，只是\emph{各维护一个}。常见三档：

\begin{itemize}
  \item \textbf{运动物体分割}（MOS）。最小方案是\emph{逐点}判别输入里所有动态像素/点（不再只判选定特征），手段仍是几何一致性（相机-模型跟踪的逐点配准残差、或地图中融合后的不一致）与/或语义分割，再辅以聚类去噪。学习式逐点动态分割（直接吃单帧或序列）能注入“典型运动物形状/运动模式”的先验，但跨域泛化是其软肋。
  \item \textbf{参数化模型}。若物体类别已知，可用\emph{参数模型}表示其形状，把每次测量当作对真实参数的一次观测、连同相机运动与背景\emph{联合}估计——这与 \cref{sec:dynamic-tracking} 的 MOT 同构，只是把“位姿”换成“位姿 $+$ 形状参数”。在捕捉人体上尤其成功（骨架模型，或稠密 SMPL 网格）。
  \item \textbf{同时跟踪与重建}。要重建\emph{任意}运动物体（无类别先验），常作刚体假设，为每个物体\emph{与}背景各建一个稠密模型，按“检测 $\to$ 对背景跟相机 $\to$ 对当前帧跟物体 $\to$ 把测量融进各自模型”四步走。物体多用 \textbf{TSDF}（\cref{sec:dense-tsdf}）表示，因其善于融合带噪深度——这正是 \cref{sec:intro-dynamic} 预告的 \emph{Co-Fusion}\cite{runz2017cofusion}：为背景与每个独立运动物各维护一个 TSDF、分别跟踪与融合。同族还有几个值得点名的代表\cite{schmid2026dynamic}：\emph{DirectTracker} 改用更\emph{稀疏的点云}表示，\emph{联合}估计物体的\emph{分组}与各自的刚体运动；而当物体跟踪误差把\emph{错位}测量融进模型、反过来拖累后续跟踪时，\emph{基于关键点}的跟踪（如 \emph{BundleTrack}）往往更鲁棒。把刚体假设进一步放松，则通向\emph{铰接体}（每段肢体一刚体 $+$ 关节约束）乃至一般可变形物体（DynamicFusion / KillingFusion，见 \cref{sec:dynamic-deformable}）；学习式的 \emph{AnyCam} 则用一个 Transformer 网络（与 \cref{ch:learning} 的 ViT 同源），从\emph{随手拍}的单目视频里\emph{联合}估出相机运动与含非刚性运动物的 4D 重建。
\end{itemize}

\begin{insight}[稠密动态 $=$ 把 \cref{ch:dense} 的“一个模型”复制成“一物一模型”]
\cref{ch:dense} 默认整个场景\emph{共用一张}稠密地图（一个 TSDF 体素场 / 一张 surfel 地图）。稠密动态唯一的结构性改动，是\emph{为背景与每个运动物体各开一份}相同的表示、各自用其测量融合。难点也随之而来：物体跟踪的误差会把\emph{错位}的测量融进物体模型，污染重建、反过来又拖累跟踪——这与 \cref{ch:dense} 里“位姿漂移会让 TSDF 重影”是同一种病，只是被分摊到了每个物体上。这条“复用表示、刚性升级为分体”的思路，在下一节会进一步推到\emph{非刚性}——当物体本身会变形，连“一物一刚体模型”都不够了。
\end{insight}

% ════════════════════════════════════════════════════════════════
\section{可变形 SLAM：当整个世界都在形变}
\label{sec:dynamic-deformable}

\paragraph{动机：没有静止背景可言。}
前两节的物体再怎么动，背后总假定有一块\emph{静止背景}供相机定位。但有些场景\emph{没有任何静止部分}——最典型的是\emph{微创手术}：唯一的传感器是单目内窥镜，看到的只有不断形变的器官\cite{schmid2026dynamic}。此时既要估相机相对\emph{持续形变的场景}的自运动，又要估场景的形变历史。这就是\emph{可变形 SLAM}（deformable SLAM）。由于缺了静止部分这个“锚”，问题严重欠约束，比静态、动态 SLAM 都更难。本节先厘清它与经典 NRSfM 的关系（\cref{sec:dynamic-nrsfm}），再讲有深度时的非刚性融合（\cref{sec:dynamic-rgbd-deform}，复用 \cref{ch:dense}），最后简述单目流水线（\cref{sec:dynamic-mono-deform}）。

\subsection{NRSfM 与可变形 SLAM：两个近亲}
\label{sec:dynamic-nrsfm}

\paragraph{非刚性运动恢复结构（NRSfM）。}
正如刚性视觉 SLAM 与 SfM 互为近亲，可变形 SLAM 也有一个经典的计算机视觉近亲——\emph{非刚性运动恢复结构}（Non-Rigid Structure from Motion, NRSfM）\cite{schmid2026dynamic}。NRSfM 用单目视频建模\emph{形变中}的三维场景。它最棘手的一点是：\textbf{每个三维点在每一帧都有\emph{不同}的三维位置}，于是地图不再是“一组固定点”，而是\emph{每个点一条三维轨迹}。单目下这严重欠约束，必须靠\emph{先验}：常用\emph{时序平滑}先验（点的轨迹随时间光滑）与\emph{空间}先验（一点的形变与其邻居相似）。NRSfM 通常\emph{批量}处理整段视频（离线，第三轴的离线端，\cref{sec:dynamic-axes}）。

\paragraph{两种约束形变的范式。}
\begin{itemize}
  \item \textbf{模板法（Shape from Template, SfT）}。假定已知物体的\emph{静止时形状}（带纹理的“模板” $\mathcal{M}_c$），用它去解释当前形变，大大简化问题。SfT 取决于\emph{形变模型}：一是\emph{解析等距模型}——假定表面上两点的\emph{测地距离守恒}（曲面可弯不可拉），适定且能实时；二是\emph{能量法}——联合最小化形变能量与重投影误差（下节自补能量定义）。
  \item \textbf{正统 NRSfM}。无 SfT 的强模板先验。早期假定\emph{正交相机}（远距物体才成立），用\emph{低维形状基}从各帧恢复三维轨迹，再加空间/时序/拓扑正则。近景序列（如内窥镜）必须用\emph{透视}相机模型，且\emph{等距}假设在此也很有效。
\end{itemize}

\paragraph{可变形 SLAM 与 NRSfM 的分野：浮动地图歧义。}
NRSfM 经典设定是\emph{相机固定}、物体作非刚性运动（含一个隐含的刚性变换 $+$ 局部形变，刚性部分从不显式分出）。可变形 SLAM 则反过来：\emph{相机在动、场景在形变}，目标是\emph{同时}估相机轨迹与场景形变历史。其核心假设是——\textbf{观测到的形变中，凡能用一个刚性变换解释的部分，都归入\emph{相机轨迹}；剩下的才是场景的真形变}\cite{schmid2026dynamic}。但“刚性 vs 非刚性”的切分本身是\emph{病态}的，称\textbf{浮动地图歧义}（Floating Map Ambiguity）：你既可以说“相机往右移、场景没动”，也可以说“相机没动、场景整体往左飘”——两者解释同一段观测。

\begin{insight}[浮动地图歧义 $=$ 尺度歧义的“非刚性版”]
\cref{thm:intro-scale} 讲过单目的\emph{尺度}歧义：把相机运动与场景大小同时缩放任意倍，图像序列不变。浮动地图歧义是它的\emph{近亲}：把“相机的刚性运动”与“场景的整体刚性漂移”互相挪用，观测也不变。两者都源于“单目 $+$ 缺锚”，也都靠\emph{加先验}消解——尺度歧义靠已知基线/IMU/学习式度量深度（\cref{ch:learning}）钉死尺度；浮动地图歧义靠对\emph{形变本身}加先验（位移要小、要光滑、要近等距——下节的能量正则）把“整体漂移”这种廉价解释压下去，逼优化把刚性部分尽量归给相机。这是“同一种病、同一类药”的又一例。
\end{insight}

\subsection{自补：非刚性重建的形变能量泛函}
\label{sec:dynamic-energy}

\paragraph{为什么这一节要“自补”。}
上文反复提到“形变能量”“ARAP 正则”“等距/Killing 正则”，Handbook 却\emph{只给名字与一句直觉、不给任何数学式}（这是其综述定位使然）。本书秉持自包含原则，\textbf{在此自行补齐这些能量的定义与正则项}——它们是经典几何处理（geometry processing）的标准结论，与 \cref{ch:dense} 的稠密表示天然咬合，故由本书独立讲授。先把舞台搭好：我们要把 \cref{ch:dense} 的\emph{刚性}配准（一个全局 $\mathbf{T}\in\SE(3)$ 对齐两片几何）升级为\emph{非刚性}配准（一个\emph{逐点变化}的形变场 $\mathcal{W}$）。

\paragraph{嵌入式形变图（ED graph）：把形变场离散化。}
若给每个点都配一个独立的 $\SE(3)$ 变换，自由度爆炸且无法约束。\emph{嵌入式形变图}\cite{sumner2007embedded} 的办法是：在物体上\emph{稀疏地}撒 $M$ 个\textbf{形变节点} $\mathbf{g}_m$（$m=1,\dots,M$），每个节点带一个\emph{局部}仿射/刚体变换 $(\mathbf{A}_m,\mathbf{t}_m)$（或 $(\mathbf{R}_m,\mathbf{t}_m)$，$\mathbf{R}_m\in\SO(3)$）与影响半径 $r_m$。任意点 $\mathbf{p}$ 的形变由其\emph{邻近节点}的变换\textbf{蒙皮插值}（skinning）得到：
\begin{equation}
  \mathcal{W}(\mathbf{p})=\sum_{m\in\mathcal{N}(\mathbf{p})} w_m(\mathbf{p})\,\big[\mathbf{R}_m(\mathbf{p}-\mathbf{g}_m)+\mathbf{g}_m+\mathbf{t}_m\big],
  \qquad
  w_m(\mathbf{p})=\frac{\hat w_m(\mathbf{p})}{\sum_{m'\in\mathcal{N}(\mathbf{p})}\hat w_{m'}(\mathbf{p})},
  \label{eq:def-skinning}
\end{equation}
其中\emph{蒙皮权重} $w_m(\mathbf{p})$ 由到节点的距离按高斯衰减并归一化，$\hat w_m(\mathbf{p})=\exp\!\big(-\|\mathbf{p}-\mathbf{g}_m\|^2/(2 r_m^2)\big)$（离节点越近、受其变换影响越大）。\cref{eq:def-skinning} 把“连续形变场”压成“$M$ 个节点变换”，自由度大降、查值只需对邻近几个节点插值——这与 \cref{sec:dense-surfel} 面元只在表面分配元素、与 \cref{ch:nlopt} 因子图只在局部连边，是同一种“稀疏化以求可解”的智慧。

\begin{definition}[as-rigid-as-possible（ARAP）形变能量]\label{def:dyn-arap}
ED 图的局部变换不能各行其是，否则物体会被撕裂。\emph{ARAP}（尽可能刚性）正则\cite{sorkine2007arap} 要求\emph{相邻节点的相对位置在形变前后尽量只差一个旋转}（即局部尽量刚性、不拉伸）。设节点 $\mathbf{g}_m$ 变形后位置为 $\mathbf{g}_m'$、其局部旋转为 $\mathbf{R}_m$，则
\begin{equation}
  E_{\mathrm{arap}}
  =\sum_{m=1}^{M}\sum_{n\in\mathcal{N}(m)} w_{mn}\,
  \big\|\,(\mathbf{g}_m'-\mathbf{g}_n')-\mathbf{R}_m(\mathbf{g}_m-\mathbf{g}_n)\,\big\|^2,
  \label{eq:def-arap}
\end{equation}
其中 $\mathcal{N}(m)$ 是节点 $m$ 的图邻居、$w_{mn}>0$ 是边权（常取余切权或 1）。$E_{\mathrm{arap}}=0$ 当且仅当每个局部邻域恰好经历一个刚体运动——即“最刚性”的形变。
\end{definition}

\begin{derivation}[为何 ARAP 能项要“减去一个旋转”]
朴素的想法是直接罚边长变化 $\big\|\|\mathbf{g}_m'-\mathbf{g}_n'\|-\|\mathbf{g}_m-\mathbf{g}_n\|\big\|$。但它\emph{允许}邻域整体自由旋转（旋转不改边长），却\emph{惩罚}正常的弯曲。ARAP 的高明在于：先给每个节点配一个\emph{最优局部旋转} $\mathbf{R}_m$，再罚“去掉这个旋转后”相对位置的残差（\cref{eq:def-arap}）。于是\emph{纯旋转}（弯曲）不受罚、唯有\emph{拉伸/剪切}才受罚——这才是“尽可能刚性”的精确含义。实际求解用\emph{交替最小化}（\cref{ch:nlopt} 的块坐标下降）：固定位置 $\{\mathbf{g}'\}$，每个 $\mathbf{R}_m$ 有\emph{闭式}解（对邻域协方差矩阵做 SVD、取最近旋转，正是 \cref{ch:vo} 3D-3D 对齐里的 Procrustes/Kabsch）；再固定 $\{\mathbf{R}_m\}$，对 $\{\mathbf{g}'\}$ 求解是一个\emph{线性}最小二乘。两步迭代，与 ICP 的“猜对应-求变换”交替（\cref{sec:lidar-scan-reg}）异曲同工。
\end{derivation}

\begin{definition}[近等距 / Killing 正则]\label{def:dyn-killing}
另一条压住形变的先验是\emph{近等距}——形变应尽量保持\emph{表面内蕴距离}（可弯不可拉、不可撕）。在连续视角下，一个\emph{瞬时}保持度量的速度场 $\mathbf{u}(\mathbf{x})$ 称\textbf{Killing 向量场}，其判据是\emph{对称化梯度为零}：$\nabla\mathbf{u}+\nabla\mathbf{u}^\top=\mathbf{0}$。真实形变不必严格等距，故\emph{近似 Killing}（Killing/近等距）正则把这一项的偏离作为惩罚：
\begin{equation}
  E_{\mathrm{kill}}=\int_{\mathcal{M}_c}\big\|\nabla\mathbf{u}(\mathbf{x})+\nabla\mathbf{u}(\mathbf{x})^\top\big\|_F^2\,\mathrm{d}\mathbf{x}
  \;\xrightarrow{\text{离散化}}\;
  \sum_{m}\sum_{n\in\mathcal{N}(m)} \big\|\,(\nabla\mathbf{u})_{mn}+(\nabla\mathbf{u})_{mn}^\top\,\big\|_F^2,
  \label{eq:def-killing}
\end{equation}
其中 $\mathbf{u}$ 是 ED 节点处形变场的瞬时速度、$(\nabla\mathbf{u})_{mn}$ 是其在边 $(m,n)$ 上的离散梯度，$\|\cdot\|_F$ 为 Frobenius 范数。\textbf{与 ARAP 的区别}：ARAP 罚“局部不刚性（拉伸/剪切）”、关心\emph{相对位置}；Killing 罚“局部不等距”、关心\emph{速度场的对称部分}，对\emph{拓扑变化}（如两只手合拢成一体）更友好——这正是 KillingFusion 能处理拓扑变化的原因。
\end{definition}

\begin{insight}[三种正则，一个目的：消解浮动地图歧义]
ED 蒙皮（\cref{eq:def-skinning}）只负责\emph{把形变场参数化}；真正\emph{约束}解的是 ARAP（\cref{eq:def-arap}）或 Killing（\cref{eq:def-killing}）这类\emph{正则项}——它们给“形变”一个“尽量小、尽量刚、尽量等距”的偏好，从而把 \cref{sec:dynamic-nrsfm} 的浮动地图歧义压下去。再加一条\emph{粘弹性}（visco-elastic）时序正则（罚形变随时间的突变，精神同 NRSfM 的时序平滑先验），就构成单目可变形 SLAM 常用的全套先验。这与本书一贯的立场一致：\emph{欠约束问题靠先验定解}——尺度歧义靠度量先验、过拟合靠鲁棒核、轨迹靠运动先验（\cref{sec:est-steam} 的 GP 运动先验），形变靠形变能量正则。
\end{insight}

\subsection{有深度时：非刚性融合（复用 \texorpdfstring{\cref{ch:dense}}{第~5~章}）}
\label{sec:dynamic-rgbd-deform}

\paragraph{把刚性 TSDF 融合升级为非刚性。}
当有 RGB-D 或双目深度时，问题\emph{过约束}、远比纯单目可解\cite{schmid2026dynamic}。这一支的奠基之作正是 \cref{sec:intro-dynamic} 预告的 \textbf{DynamicFusion}\cite{newcombe2015dynamicfusion}。它的骨架\emph{完全建立在 \cref{ch:dense} 之上}，只把一处“刚性”换成“非刚性”：

\begin{enumerate}
  \item \textbf{标准模型 $\mathcal{M}_c$}：维护一个\emph{静止时形状}的 TSDF 体素场（就是 \cref{sec:dense-tsdf} 的 TSDF，只是约定它表示“canonical”姿态）；
  \item \textbf{形变场 $\mathcal{W}$}：用 ED 图（\cref{eq:def-skinning}）把标准模型\emph{非刚性地}变形到\emph{当前帧}的姿态，配 ARAP 正则（\cref{eq:def-arap}）保持局部刚性；
  \item \textbf{配准（求 $\mathcal{W}$）}：最小化“变形后的标准模型”与“当前深度观测”的数据项 $+$ ARAP 正则项
  \begin{equation}
    E_{\mathrm{def}}=\underbrace{\sum_{\text{深度点}}\big\|\,\mathrm{psdf}\big(\mathcal{W}(\mathbf{v})\big)\,\big\|^2}_{E_{\mathrm{data}}：\text{对齐当前深度}}
    +\ \lambda\,\underbrace{E_{\mathrm{arap}}}_{E_{\mathrm{reg}}：\text{局部尽量刚性}},
    \label{eq:def-fusion-energy}
  \end{equation}
  其中数据项把标准模型表面顶点 $\mathbf{v}$ 经形变场 $\mathcal{W}$ 送到当前帧、取其\emph{投影式符号距离} $\mathrm{psdf}$（即 \cref{eq:dense-tsdf-proj} 的投影式 SDF，沿视线测到最近表面的符号距离），$\lambda$ 权衡数据与正则；
  \item \textbf{融合（更新 $\mathcal{M}_c$）}：求得形变场后，把当前深度\emph{反向}经 $\mathcal{W}^{-1}$ 拉回标准模型，按 \cref{ch:dense} 的\emph{加权运行平均}（\cref{eq:dense-tsdf-fuse}）融进 TSDF——\textbf{这一步与静态 KinectFusion 的融合\emph{逐字相同}}，唯一区别是测量先被“拉回标准姿态”。
\end{enumerate}

\begin{insight}[非刚性融合 $=$ \cref{ch:dense} 的 TSDF $+$ 一层形变场]
盯住 \cref{eq:def-fusion-energy} 与 \cref{ch:dense} 的对照：静态 KinectFusion 用\emph{一个}全局位姿 $\mathbf{T}\in\SE(3)$ 做 frame-to-model 配准（\cref{sec:dense-tsdf}），DynamicFusion 把这个\emph{单一刚性变换}换成\emph{逐点形变场} $\mathcal{W}$（\cref{eq:def-skinning}）、并加 ARAP 正则约束它——\emph{除此之外，TSDF 的截断、投影式 SDF、加权平均融合、光线投射提面，统统照搬}。\cref{eq:dense-tsdf-fuse} 那个把累积权重封顶在 $W_\eta$ 的“移动平均”当初就埋了“支持动态场景”的伏笔，这里正式兑现。\cref{sec:dense-surfel} 还讲过 ElasticFusion 的“无位姿图、直接对面元地图做非刚性形变（deformation graph）”——那已是非刚性融合的雏形；SurfelWarp 进一步用 surfel（\cref{sec:dense-surfel}）替代 SDF 以改善可扩展性。可见 \cref{ch:dense} 的整套稠密机器，\emph{加一层形变场}就长出了处理可变形场景的能力。
\end{insight}

\paragraph{谱系与医学应用。}
沿此线，VolumeDeform 加入光度信息改善结果；\emph{KillingFusion}\cite{slavcheva2017killingfusion} 把 ARAP 换成\emph{Killing}正则（\cref{eq:def-killing}），从而能处理\emph{拓扑变化}（如双手合拢）；SurfelWarp 改用面元提升可扩展性。在医学领域，MIS-SLAM 为双目内窥镜做基于 ARAP $+$ ED 的可变形 SLAM；EMDQ 则用“期望最大化 $+$ 对偶四元数”在医学立体下跟踪相机并估组织形变。这些系统\emph{共享}上述“TSDF/surfel $+$ 形变图 $+$ 能量正则”的骨架，差别只在正则项与表示的取舍——故本书提炼共性、不逐一复述。

\subsection{单目可变形 SLAM 的流水线（选读）}
\label{sec:dynamic-mono-deform}

\paragraph{跟踪 $+$ 建图，照搬视觉 SLAM 的双线程。}
纯单目可变形 SLAM（如 DefSLAM、NR-SLAM）模仿视觉 SLAM 的\emph{帧率跟踪 $+$ 关键帧率建图}双线程（\cref{ch:vo}），只是把每一环“非刚性化”\cite{schmid2026dynamic}：
\begin{itemize}
  \item \textbf{跟踪}（帧率）：类似 SfT——假定当前探索区域的\emph{静止时形状}可得，用\emph{物理弹性/粘弹性 $+$ ARAP $+$ 时序}正则的形变能量，估当前帧形变与相机位姿。
  \item \textbf{建图}（关键帧率）：吃点轨迹，在关键帧率\emph{重算模板}。DefSLAM 用等距 NRSfM 从若干共视关键帧算出每关键帧一个模板；NR-SLAM 引入\emph{动态可变形图}应对任意拓扑、并把粘弹性形变模型嵌进核心的\emph{可变形 BA}（即把 \cref{eq:def-fusion-energy} 式能量放进 \cref{ch:nlopt} 的束调整里联合优化形变与位姿）。
  \item \textbf{特征匹配}：医学近景里 ORB（\cref{ch:vo}）表现不佳，\emph{Lucas--Kanade 光流}（\cref{ch:vo} 的直接法/光流，对仿射光照变化不变）更稳；重定位仍用 ORB $+$ 词袋（\cref{ch:loop}）。
\end{itemize}

\begin{pitfall}[陷阱：以为“可变形 SLAM 就是把刚性 SLAM 的 $\mathbf{T}$ 换成形变场”就完了]
\textbf{现象}：直接把 \cref{ch:vo} 的重投影 BA 里相机位姿换成逐点形变场，结果优化\emph{发散}或收敛到“场景整体乱飘”。\textbf{根本原因}：少了\emph{两样东西}——(1) 没加形变正则（\cref{sec:dynamic-energy} 的 ARAP/Killing/时序），浮动地图歧义（\cref{sec:dynamic-nrsfm}）没被消解；(2) 没有\emph{初始化}与\emph{模板}，单目可变形从零启动比刚性更难（DefSLAM 靠等距 NRSfM 起步、NR-SLAM 靠刚性 SfM 给初值再切非刚性）。\textbf{正确做法}：形变场 $+$ \emph{足量先验} $+$ 可靠初始化，三者缺一不可。\textbf{自检}：你的可变形优化目标里，正则项与数据项的权重 $\lambda$（\cref{eq:def-fusion-energy}）调过吗？去掉正则会怎样？
\end{pitfall}

% ════════════════════════════════════════════════════════════════
\section{长期与终身 SLAM：地图随时间变化}
\label{sec:dynamic-longterm}

\paragraph{动机：变化发生在“你没看的时候”。}
前面处理\emph{短期}动态（视野内）。本节转向\emph{长期}动态（\cref{def:dyn-short-long}）——变化发生在机器人\emph{视野之外}，它只看到累积的、突兀的结果。长期动态给 SLAM 带来几重新挑战\cite{schmid2026dynamic}：其一，变化\emph{突兀且可任意大}（机器人越久不看某处，那处越可能面目全非），数据关联从“跟踪”退化为\emph{全局重识别与匹配}，且\emph{感知混叠}（\cref{sec:intro-need}）加剧——一个独特路标可能因移动而在多处被看到；其二，许多在短期可忽略的效应（物体增删移、结构改建、光照/季节/纹理变化）在长期变得显著；其三，长期常伴随大视角变化与传感器退化。本节给一套关于“地图随时间变化”的概念框架与可识别性论证，并把它与 \cref{ch:loop} 的回环、\cref{sec:dynamic-energy} 的形变接口对齐。

\subsection{终身定位：在变化中认出同一处地方}
\label{sec:dynamic-lifelong-loc}

\paragraph{与回环同源。}
终身 SLAM 的核心能力之一，是\emph{在上述挑战下仍能准确定位}。这与 \cref{ch:loop} 的\emph{地点识别/回环检测}概念上高度一致——都是“尽管变了，仍认出是同一处”。几条主流做法：
\begin{itemize}
  \item \textbf{并行/汇总地图（experiences/sessions）}。把每次造访称\emph{经历}（experience）或\emph{会话}（session），每次经历内变化小到可跟踪。若当前经历能相对\emph{以往经历}定位，就说明地图够用、无需加点；定位弱或失败才把新数据\emph{添}进地图——地图复杂度因而\emph{自适应}于外观变化的程度。可加\emph{跨经历约束}（$\mathcal{E}_A\!\to\!\mathcal{E}_B\!\to\!\mathcal{E}_C$）提升一致性，并\emph{离线}跑全局 BA 汇总、只保留最利于未来定位的特征。
  \item \textbf{外观不变表示}。不为每种外观各存一份特征，而是训练能\emph{跨外观匹配}的检测/描述/匹配（深度学习之力，\cref{ch:learning,ch:loop}）。一种特殊的外观不变特征是\emph{语义}：“一把椅子”在不同光照、不同视角下\emph{仍是椅子}，故物体及其布局可跨大视觉变化用于定位（如空地视角与地面视角间）。
  \item \textbf{记忆、扩展与边缘化}。终身运行的机器人不断收集信息、地图无限膨胀。对策有三：\emph{遗忘}（给地图点一个随时间衰减的\emph{持久概率}，久未观测者降权直至剪除）；\emph{记忆管理}（如 RTAB-Map 在短期/工作/长期记忆间调度节点，保证当前处理的节点数足够实时）；\emph{边缘化}（见下文陷阱）。
\end{itemize}

\begin{pitfall}[陷阱：边缘化把错误“烤进”地图（baking-in）]
\textbf{现象}：滑窗 SLAM 为控制规模，把旧位姿/旧因子\emph{边缘化}（\cref{eq:info-marginal} 的舒尔补，用一个概括分布替换被移除的节点）。但一旦某条\emph{回环}的内/外点判定还没定准就被边缘化，它的当前赋值会被\textbf{冻结进}那个概括节点、\emph{再不能更改}——即便后来出现了否定它的证据。\textbf{根本原因}：边缘化\emph{保留信息但牺牲灵活性}——被边缘化节点的\emph{拓扑不可再改}（\cref{ch:nlopt} 已点明这一代价）。\textbf{正确做法}：对\emph{尚有歧义}的因子（尤其回环）\emph{延迟}边缘化，待证据充分再固化；何时边缘化是一个重要的开放问题。\textbf{自检}：你边缘化掉的因子里，有没有“将来可能被推翻”的回环或动态判定？
\end{pitfall}

\subsection{地图清理与变化检测：无证据 vs 有缺席证据}
\label{sec:dynamic-change-detect}

\paragraph{一个看似哲学、实则要命的区分。}
即便有遗忘机制，检测式方法仍会\emph{累积过时观测}（如早已移走的物体）。根子在于一个不对称：\textbf{探测器“看到”是某物\emph{存在}的证据；但“没看到”\emph{不}等于它\emph{不存在}}\cite{schmid2026dynamic}。这就是\textbf{“无证据”（absence of evidence）vs“有缺席证据”（evidence of absence）}问题。

\begin{insight}[absence of evidence $\neq$ evidence of absence]
机器人某次造访某房间，在角落看到一把椅子（正探测）。下次重访\emph{没看到}它，有三种可能：(1) 它\emph{没往那个角落看}——椅子可能还在（\emph{无证据}，absence of evidence）；(2) 它\emph{看了那个角落、确认椅子不见了}（\emph{有缺席证据}，evidence of absence，即真正的\emph{负观测}）；(3) 它\emph{看了、但因光照变化没认出}椅子——也会误判“椅子没了”。只有第 (2) 种才是\emph{变化}。\textbf{要区分这三者，地图表示必须显式区分“未观测”与“观测到为空”}——这正是 \cref{sec:dense-octomap} 的占据栅格/TSDF 把空间分成\emph{占据/空闲/未知}三态、而非简单二态的深层用意：\emph{“未知” $=$ 无证据，“空闲” $=$ 有缺席证据}。一张只存“看到的表面点”的稠密点云（\cref{sec:dense-rgbd-pc}）做不到这一点（它没有“观测到为空”的概念），这也是为什么\emph{变化检测偏爱体素/占据表示}。
\end{insight}

\paragraph{地图清理与几何/语义变化检测。}
\emph{地图清理}（map cleaning）旨在从地图中移除动态与虚假测量、只留高质量静态特征。离线处理时，先用鲁棒 BA 估全局一致位姿（动态点当外点），再重访所有测量查一致性：要么把它们融进一个全局\emph{体素地图}（占据/TSDF）逐格查一致（全面但费内存），要么用\emph{可见性}方法只取邻近测量子集查（省）。\emph{变化检测}（change detection）则多在\emph{多会话}设定下：每会话各重建一张内部一致的（已清理的）地图，再在\emph{场景层}找两次造访间的变化。手段有二，常联用：
\begin{itemize}
  \item \textbf{几何变化检测}：用\emph{体素差分}（scene differencing）比对两张占据/TSDF 地图——\emph{曾观测为空、现出现表面}者为新增；\emph{曾有表面、现观测为空}者为消失。要先把各会话\emph{联合优化}对齐（不能只刚性对齐，须用跨会话约束 $+$ BA 抵消会话内漂移）。
  \item \textbf{语义变化检测}：语义标签\emph{外观不变}，是跨条件比对的强抽象；且能捕捉\emph{非几何}变化（桌上撤走一张纸、墙刷成别的颜色）。更关键的是——\emph{物体常作为“一个连贯单元”整体变化}，这与闭集语义类别高度重合（杯子、椅子整体移动）。
\end{itemize}

\begin{pitfall}[陷阱：在部分观测上做变化检测，留下“半张桌子”]
\textbf{现象}：机器人重访发现桌子不见了，但视野只覆盖了桌子一半，于是只删掉\emph{看到的那半}，地图里飘着\emph{半张桌子}——对人而言荒谬（桌子总是整体移动）。\textbf{根本原因}：\emph{语义一致性}（semantic consistency）问题——场景按\emph{连贯单元}变化，机器人却只拿到\emph{部分}观测。\textbf{正确做法}：把变化检测提到\emph{物体/语义类别}的层级（“整把椅子”而非“若干表面点”），由\emph{构造}保证语义一致；或每次从零重建（保一致但永远建不出完整模型）。\textbf{自检}：你的变化单元是“一团点”还是“一个物体”？
\end{pitfall}

\subsection{物体级理解与实例重定位：从变化中认识物体}
\label{sec:dynamic-object-understanding}

\paragraph{从变化反推物体（Object Detection from Changes）。}
变化与语义物体之间有一层\emph{双向}的内在关联，两个方向都可利用\cite{schmid2026dynamic}。\cref{sec:dynamic-change-detect} 已用过一向——\emph{语义改进变化检测}（物体作连贯单元整体变化）；反过来，“某处\emph{变了}”本身也是\emph{物体检测}的线索：在稠密重建里检出所有\emph{运动表面}、再\emph{空间聚类}，便能识别出可能的物体。这种几何与语义的\emph{互补}在两种情形尤为有用：根本\emph{没有}语义检测器时，或用来\emph{补全}语义分割漏掉的物体。但要警惕——\emph{纯几何}的物体分割对\emph{部分观测}敏感：两个接触的物体一起被移走，易\emph{欠分割}成一个；被遮挡而只看到一部分的物体，易\emph{过分割}成数块（这与 \cref{sec:dynamic-change-detect} 的“半张桌子”陷阱同根：场景按连贯单元变化，观测却是局部的）。

\paragraph{物体实例重定位：与回环同源，但更难。}
要建立\emph{物体中心}的\emph{时序}理解（而非只估当前或静态状态），就需\emph{物体实例重定位}（object instance re-localization）：在长期动态观测里\emph{跨时间跟踪同一个物体}。概念上，这个问题及其多数解法又一次与\emph{地点识别 / 回环检测}（\cref{ch:loop}）高度相似——都是“尽管变了、仍认出是同一个”。但有三个额外难点使它比回环更棘手\cite{schmid2026dynamic}：
\begin{itemize}
  \item \textbf{突兀的大位移}。长期动态的突兀本性（\cref{def:dyn-short-long}）让物体在两次观测间\emph{原则上可任意远地}移动或改变。物体\emph{上次被观测的状态}（结合流逝时间与假定的运动模型）\emph{有时}是关联的有用线索，但不总成立。
  \item \textbf{物体小、特征少}。单个物体远小于整个场景，几何与纹理的\emph{多样性更少}，能唯一刻画它的特征也更少；小物体若没被\emph{近距离}观测，可用测量更少、分辨率更低，雪上加霜。
  \item \textbf{相同实例的感知混叠}。环境里常有多个\emph{相同或相似}类型的实例，\emph{感知混叠}（\cref{sec:intro-need}）成为主要挑战，并牵出一个看似哲学的问题：物体“是一个唯一实例”究竟意味着什么、是否该把它们逐个跟踪。
\end{itemize}

\begin{insight}[感知混叠与“该不该逐个跟踪”：唯一性决定多假设]
设想一间研讨室里有几十上百把\emph{一模一样}的椅子：多半分不清哪把是哪把、各自如何被移动过。但\emph{这对多数任务也许无关紧要}——不如把所有椅子的观测\emph{融成一个有代表性的模型、再复制多份}。反例是你\emph{自己}的咖啡杯：办公楼里也许有无数相同的杯子，可人们\emph{很在意}追踪\emph{自己那只}。可见“相同实例该不该逐个跟踪”由\emph{任务}定、没有先验答案——这与 \cref{ch:openworld} “物体粒度由任务决定”的信息瓶颈洞见同源。应对感知混叠的一条务实原则是\textbf{考虑解的唯一性}：高度独特的物体若有\emph{唯一}匹配，可径直接受；有\emph{多个}相似候选的物体，则改用\emph{多假设}方法、或给出恰当的\emph{不确定度度量}，把判断\emph{推迟}到更多证据到来——这正与 \cref{ch:loop} 对付错误回环、\cref{sec:dynamic-lifelong-loc} 对“歧义因子延迟边缘化”所共享的智慧。
\end{insight}

\paragraph{闭集 vs 开集重定位。}
物体实例重定位可在\emph{闭集}或\emph{开集}两种设定下考虑\cite{schmid2026dynamic}。\emph{闭集}假定空间里的物体集\emph{固定}，于是重定位本质上是一个\emph{最优匹配}问题（把当前观测指派到已知物体集）。\emph{开集}（更一般）允许物体被\emph{增删}，于是多出一问——“物体观测\emph{相似到什么程度}才算同一个？”因为总存在一个替代解释：场景里\emph{新出现}了一个相似但不同的物体。这与 \cref{ch:openworld} 从闭集语义到开集语义的跨越\emph{同构}：一旦容许“没见过的新物体”，匹配就不再是封闭集合上的指派，而要为“这是新东西”留出可能。

\paragraph{描述子族：从 FPFH/SHOT 到学习式、从多关键点到单描述子。}
主流解法聚焦于\emph{描述子的提取与匹配}\cite{schmid2026dynamic}，可分两族。

\emph{第一族——每物体一组关键点 $+$ 描述子}：在每个物体上提一组\emph{关键点}并为各点算描述子，直接用于找对应、解 6 自由度（6-DoF）变换。最自然的第一步，是直接\emph{复用 SLAM 前端的视觉特征}（\cref{ch:vo}）作关键点。或者用\emph{三维}关键点与描述子，如 \textbf{FPFH}（Fast Point Feature Histograms，快速点特征直方图）与 \textbf{SHOT}（Signature of Histograms of OrienTations，方向直方图签名）——二者在\emph{变化物体重定位}里用得很多。其优点有二：可从一次会话的\emph{稠密三维重建}里提取、且\emph{对外观变化不敏感}（不像光度特征那样随光照、季节漂移）。代价是：它们只刻画\emph{一小块局部几何}，故描述力较弱、常需进一步的\emph{几何验证}（如 RANSAC 式一致性核验，\cref{ch:vo}）去伪。近来\emph{深度学习}式的局部形状（有时含外观）描述子性能最佳——因为可\emph{训练}它去强调一个形状或物体里\emph{最具区分力}的部分。

\emph{第二族——每物体一个描述子}：不提多个关键点，而为每个物体实例算\emph{单个}描述子\cite{schmid2026dynamic}。好处是省去维护一个庞大关键点库、并避免对它\emph{穷举匹配}的高昂计算。其\emph{关键}在于：要把网络\emph{构造或训练}成 $\SE(3)$-\emph{不变}或 $\SE(3)$-\emph{等变}的——前者让描述子\emph{跨视角}仍能匹配，后者更进一步，允许直接\emph{配准}（从匹配上的两个物体描述子里读出它们之间的 $\SE(3)$ 变换）。这正是 \cref{sec:dynamic-change-aware} 那条洞见所点的“用网络从部分观测学完整形状潜码”的描述子家族；它们如何嵌进\emph{在线} change-aware 因子图做鲁棒数据关联，见 \cref{sec:dynamic-change-aware}。

\begin{note}[$\SE(3)$-不变 vs 等变：一字之差，能否配准]
设 $\mathbf{d}(\cdot)$ 为物体的描述子、$\mathbf{T}\!\cdot\! o$ 表示把物体 $o$ 经 $\mathbf{T}\in\SE(3)$ 摆到新位姿。\textbf{不变}（invariant）指 $\mathbf{d}(\mathbf{T}\!\cdot\! o)=\mathbf{d}(o)$——描述子对位姿\emph{不敏感}，只够判“是不是同一物体”。\textbf{等变}（equivariant）指 $\mathbf{d}(\mathbf{T}\!\cdot\! o)=\mathbf{T}\!\circ\!\mathbf{d}(o)$（$\circ$ 为 $\SE(3)$ 在描述子上的群作用），描述子随物体\emph{协同}变换，故把\emph{姿态}也编码进来，能从两个匹配描述子\emph{解出}相对位姿。一句话：\emph{不变够识别、等变才能配准}——这与本书 \cref{ch:loop} 区分“地点识别（判同处）”与“重定位（求位姿）”是同一层意思。
\end{note}

\noindent 物体实例重定位的代表性基准是 \textbf{RIO} 数据集：给定一处\emph{变化室内环境}的\emph{初始扫描}与之后的\emph{重扫}，目标是认出哪些物体实例相互对应、并估出每个物体到其\emph{新位置}的 6-DoF 位姿\cite{schmid2026dynamic}。最后值得一提的是这一能力的\emph{增益}：一旦变化物体的众多\emph{部分观测}能被重定位与配准，不仅给出每个物体随时间演化的\emph{历史}，还能把这些部分观测\emph{累积}成越来越\emph{完整}的物体模型与越来越丰富的场景描述——这又呼应 \cref{ch:openworld} 用观测增量精化物体表示的思路。

\subsection{change-aware SLAM 与统一框架}
\label{sec:dynamic-change-aware}

\paragraph{在线检测长期变化：局部一致性是命门。}
多会话只是长期动态的一个特例。更一般地，变化可能在\emph{任意}两次观测间发生（机器人去厨房一趟，客厅就变了，无论隔一周还是一分钟）。\emph{change-aware SLAM} 要在\emph{在线}运行中检测并表示这类变化。其最深的难点，是\emph{变化检测与定位互相纠缠}\cite{schmid2026dynamic}：若位姿偏了一米，先前建的物体表面在当前（偏移的）视角下会\emph{显得缺席}，明明没变也像变了。考虑两个极端解就看清了这层纠缠——全判“没变”（即便真有大变），配准会被真假测量\emph{平均}得不准；全判“变了、now absent”，当前开环估计\emph{平凡最优}，但重建残缺、不一致。

\begin{insight}[change-aware SLAM 的可识别性原理：局部一致性（local consistency）]
如何把\emph{噪声}与\emph{真变化}分开？核心思想是用\textbf{局部一致性}建立一段段\emph{保证无变化}的观测序列\cite{schmid2026dynamic}。\textbf{底线情形是逐帧跟踪}（frame-to-frame）：直觉上——\emph{机器人正盯着看的时候，场景不可能在它眼皮底下悄悄变}。在“变化率低 $+$ 短时里程计误差小”的合理假设下，这可放宽到一个\emph{滑动}或\emph{主动}窗口。于是 change-aware SLAM 的可识别性论证可凝练为：\textbf{“看着的时候不会变”给出局部不变的锚段，跨锚段的不一致才被判为真变化}——这与 \cref{ch:loop} 用“几何验证”筛真假回环、与 \cref{sec:dynamic-removal} 用“多视一致性”判动态外点，是\emph{同一种}“以局部确定性切割全局歧义”的思路。也正因如此，多数 change-aware 方法\emph{先把测量局部汇总成物体级表示}再推理——物体既是变化的连贯单元，又天然保证语义一致。
\end{insight}

\paragraph{统一短期与长期：一张图容纳多种动态。}
局部一致性的妙处在于：它能把\emph{短期动态物体跟踪}（\cref{sec:dynamic-tracking}）与\emph{长期变化检测}缝在一起，同时捕捉多种动态模式。代表性提法本书只点其机制、不复述管线\cite{schmid2026dynamic}：
\begin{itemize}
  \item \emph{Changing-SLAM} 在 ORB-SLAM 之上，按语义掩膜把特征点分到物体，用卡尔曼滤波\emph{逐帧跟踪}短期运动；\emph{未被跟踪}的点（长期变化）则\emph{全局}关联到邻近同类物体；每物体估一个持久分，久未探测者遗忘。
  \item \emph{Panoptic Multi-TSDF} 给定全局一致位姿，只更新\emph{能逐帧跟踪}且语义匹配的\emph{活跃子图}，用活跃-非活跃子图\emph{差分}高效检测变化，按一致/冲突\emph{合并或丢弃}旧子图——由“整子图删除”\emph{构造性地}保证语义一致。
  \item \emph{POV-SLAM} 用\emph{期望最大化}（EM，\cref{ch:nlopt} 的交替优化）在“给定关联优化位姿/物体”与“给定布局优化关联/持久概率”间迭代，持续精化定位与变化检测。
  \item \emph{Khronos} 做稠密时空（4D）度量-语义 SLAM：用鲁棒因子图优化（\emph{逐步非凸} GNC，\cref{ch:nlopt}）定内点关联，把\emph{背景形变}当额外因子辅助空间优化，回环后再做一次\emph{可变形}几何变化检测来消解“无证据 vs 缺席证据”。
\end{itemize}

\begin{insight}[长期 change-aware SLAM 仍是母方程 $+$ 物体级因子]
盯住上面这些系统的\emph{机制}：POV-SLAM 的 EM、Khronos 的鲁棒因子图，\emph{都在解一张以物体观测为路标的因子图}——变量是相机位姿 $+$ 物体位姿/持久性，因子是物体级的关联约束（含回环、含背景形变）。这与 \cref{sec:dynamic-tracking} 把运动物体增广进因子图\emph{同构}，只是这里物体的“运动”是长期的、突兀的，且多了“持久性/存在与否”这一维。换言之——\textbf{从短期跟踪到长期变化检测，一以贯之的还是本书的 MAP 母方程（\cref{eq:intro-nls}），只是因子的语义从“看到了哪”扩展到“它还在不在、跟谁是同一个”}。物体重定位所需的\emph{$\SE(3)$-不变/等变描述子}（用网络从部分观测学出完整形状的潜码，\cref{ch:learning}），则是这张图里\emph{鲁棒数据关联}的现代答案，与 \cref{ch:loop} 的地点识别描述子一脉相承。
\end{insight}

\paragraph{再往前：时序场景理解（展望）。}
最后一步是超越“重建变化”、走向\emph{预测变化}——\emph{时序场景理解}。简述三族：\emph{动态地图}（Maps of Dynamics）从轨迹估场景的\emph{典型运动模式}（如单行道主流向、人群早晚潮汐）；\emph{周期事件}用\emph{傅里叶频率图}（\cref{ch:imu} 的频域分析同源）把场景状态当周期过程的观测、只存主频，从而\emph{预测}未来状态、改善定位与规划；\emph{学习式方法}用场景图（scene graph）作物体与关系的紧凑表示去预测演化，能融语义先验（“餐桌中午前后会变”）。这条线把 SLAM 推向“世界模型”——把过去、现在、未来的场景演化统一起来，正是 \cref{ch:openworld} 开放世界空间 AI 与 \cref{ch:epilogue} 结语要接续的方向。

% ════════════════════════════════════════════════════════════════
\section{常见误解}
\label{sec:dynamic-myths}

\begin{pitfall}[误解一：动态 SLAM 是与经典 SLAM 无关的另一套方法]
本章反复证伪了这一点。动态物体 SLAM 只是把运动物体的位姿/速度\emph{增广}进状态、把运动模型写成因子（\cref{eq:dyn-map}），令物体位姿 $\mathbf{T}_{wi}^{k}\equiv\mathbf{I}$ 就退回静态 SLAM；可变形融合只是把 \cref{ch:dense} 的刚性配准换成形变场（\cref{eq:def-fusion-energy}）、TSDF 融合照搬；长期 change-aware 仍是物体级因子图。判断任何“动态/可变形 SLAM”的第一问永远是：\emph{它在经典框架上放松了哪条静态假设、新增了哪类变量与因子？}
\end{pitfall}

\begin{pitfall}[误解二：短期/长期动态按绝对时间长短区分]
错。\cref{def:dyn-short-long} 说清了：短期/长期取决于\emph{场景变化率与观测率之比}，是\emph{观测者}的属性——同一段运动，对盯着看的机器人是短期（连续轨迹），对看一眼就走的机器人是长期（瞬移般跳变）。判据常简化为“现在是不是在视野里看着它动”。把它当“几秒 vs 几天”，会在多会话/多机器人场景里完全错判。
\end{pitfall}

\begin{pitfall}[误解三：语义分割出“可能动的类别”就能解决动态]
不够。语义只回答“这类东西\emph{会不会}动”，不回答“它\emph{此刻}动没动”（停着的车是潜在动态、却是好路标）。必须叠加\emph{当前帧的多视几何一致性}检验（\cref{sec:dynamic-removal}）。把“语义可动”直接当“动态外点”剔掉，会误删大量静止的潜在动态物、反伤定位（DynaSLAM 用“语义 $+$ 几何”双保险正为此）。
\end{pitfall}

\begin{pitfall}[误解四：可变形 SLAM 就是把刚性 SLAM 的 $\mathbf{T}$ 换成形变场]
远不止。少了\emph{形变正则}（ARAP/Killing/时序，\cref{sec:dynamic-energy}），浮动地图歧义（\cref{sec:dynamic-nrsfm}）不会自行消解，优化会让场景“整体乱飘”；少了\emph{模板与初始化}，单目可变形从零启动比刚性更难。形变场 $+$ 足量先验 $+$ 可靠初始化，三者缺一不可。
\end{pitfall}

\begin{pitfall}[误解五：“没看到某物”就等于“它消失了”]
这是 absence of evidence 与 evidence of absence 的混淆（\cref{sec:dynamic-change-detect}）。“没看到”可能是\emph{没往那看}、\emph{看了没认出}、或\emph{真的没了}——只有第三种是变化。要区分，地图必须显式区分“未观测”与“观测到为空”，这正是占据/TSDF 用\emph{占据/空闲/未知}三态（\cref{sec:dense-octomap}）、而纯点云做不到的根由。
\end{pitfall}

\begin{pitfall}[误解六：边缘化只是“省内存的无损压缩”]
边缘化\emph{保信息但损灵活性}：被边缘化节点的拓扑\emph{不可再改}，尚有歧义的回环/动态判定一旦被边缘化就\emph{冻结}（baking-in，\cref{sec:dynamic-lifelong-loc}），未来证据也翻不了案。对长期运行的机器人，\emph{何时}边缘化、对哪些因子\emph{延迟}边缘化，是要紧的设计抉择，而非纯粹的工程优化。
\end{pitfall}

% ════════════════════════════════════════════════════════════════
\section*{本章小结}

本章是全书“放松世界模型假设”的一章，承 \cref{ch:learning}（把网络接进几何）之后，把经典 SLAM 的\emph{静态刚体世界}假设逐层松开。一条主线贯穿始终：\textbf{动态与可变形 SLAM 不推翻经典框架，而是把它\emph{推广}——往因子图里增广运动物体的位姿/速度、把运动写成因子（动态物体 SLAM），把 \cref{ch:dense} 的刚性配准换成带能量正则的形变场（可变形 SLAM），把“看到了哪”的因子扩展到“它还在不在”（长期 change-aware SLAM）}。三条松绑：\emph{短期动态}（视野内运动，剔除或增广跟踪，\cref{eq:dyn-map}）$\to$ \emph{可变形}（无静止背景，ED 图 $+$ ARAP/Killing 能量 $+$ 非刚性 TSDF 融合）$\to$ \emph{长期动态}（视野外变化，局部一致性 $+$ 占据三态 $+$ 物体级因子图）。无论松到哪一步，“优化即推断、几何融证据”这条主干始终不倒——这正是它与经典各章\emph{浑然一体}的根由。

\begin{table}[H]\centering\small
\caption{符号表}
\begin{tabular}{lll}
\toprule
符号 & 含义 & 首次出现 \\
\midrule
$\mathbf{T}_{wi}\in\SE(3)$ & 相机位姿（$i$ 系到世界系；本书下标式）& \cref{eq:dyn-reproj} \\
$\mathbf{T}_{wi}^{k}\in\SE(3)$ & 第 $k$ 个运动物体在 $i$ 时刻的位姿 & \cref{eq:dyn-reproj} \\
$\mathbf{v}_i^{k},\boldsymbol\omega_i^{k}$ & 物体 $k$ 在 $i$ 时刻的线/角速度（物体系）& \cref{eq:dyn-vcte} \\
$\mathbf{p}_j^{k},\tilde{\mathbf{p}}_j^{k}$ & 物体 $k$ 上点 $j$ 在物体系坐标（齐次）& \cref{eq:dyn-reproj} \\
$\mathbf{y}_l$ & 静态背景路标（世界系）& \cref{eq:dyn-map} \\
$\Delta\mathbf{T}^{k}_{i\to i+1}$ & 恒速模型下物体位姿增量 & \cref{eq:dyn-increment} \\
$h^{-1}$ & 齐次到欧氏的去齐次算子 & \cref{eq:dyn-couple} \\
$\mathcal{M}_c$ & 标准模型（canonical，静止时形状）& \cref{sec:dynamic-rgbd-deform} \\
$\mathcal{W}:\mathbb{R}^3\!\to\!\SE(3)$ & 形变场（把标准模型送到当前帧）& \cref{eq:def-skinning} \\
$\mathbf{g}_m,(\mathbf{R}_m,\mathbf{t}_m),r_m$ & ED 图节点 / 局部变换 / 影响半径 & \cref{eq:def-skinning} \\
$w_m(\mathbf{p})$ & 节点 $m$ 对点 $\mathbf{p}$ 的蒙皮权重 & \cref{eq:def-skinning} \\
$\mathcal{N}(i),\mathcal{N}(m)$ & 点/节点的邻居集 & \cref{eq:def-arap} \\
$E_{\mathrm{def}},E_{\mathrm{arap}},E_{\mathrm{kill}}$ & 形变能量 / ARAP 正则 / Killing 正则 & \cref{eq:def-fusion-energy,eq:def-arap,eq:def-killing} \\
$r_{\mathrm{change}}/r_{\mathrm{obs}}$ & 变化率与观测率之比（定短/长期）& \cref{def:dyn-short-long} \\
\bottomrule
\end{tabular}
\end{table}

\begin{table}[H]\centering\small
\caption{定理 / 公式速查表}
\begin{tabular}{lll}
\toprule
公式 / 概念 & 一句话说明 & 对应 \\
\midrule
短/长期定义 \cref{def:dyn-short-long} & 由变化率/观测率之比定，非绝对时长 & \cref{sec:dynamic-axes} \\
改进重投影 \cref{eq:dyn-reproj} & 比静态多插一个物体位姿 $\mathbf{T}_{wi}^{k}$ & \cref{sec:dynamic-tracking} \\
恒速残差 \cref{eq:dyn-vcte} & 相邻物体速度之差（同里程计型）& \cref{sec:dynamic-tracking} \\
位姿增量 \cref{eq:dyn-increment} & 匀速螺旋的 $\SE(3)$ 一阶积分 & \cref{sec:dynamic-tracking} \\
速度-位姿耦合 \cref{eq:dyn-couple} & “按位姿摆”$=$“按运动模型摆” & \cref{sec:dynamic-tracking} \\
动态物体 MAP \cref{eq:dyn-map} & 母方程 $+$ 三类动态新因子 & \cref{fig:dyn-factorgraph} \\
ED 蒙皮 \cref{eq:def-skinning} & 稀疏节点变换插值出形变场 & \cref{sec:dynamic-energy} \\
ARAP 能量 \cref{eq:def-arap} & 罚“去掉旋转后”的相对位移（不拉伸）& \cref{def:dyn-arap} \\
Killing 能量 \cref{eq:def-killing} & 罚速度场对称部分（近等距，容拓扑变）& \cref{def:dyn-killing} \\
非刚性融合 \cref{eq:def-fusion-energy} & \cref{ch:dense} TSDF $+$ 形变场 $+$ ARAP & \cref{sec:dynamic-rgbd-deform} \\
局部一致性 & 看着的时候不会变 $\Rightarrow$ 切割全局歧义 & \cref{sec:dynamic-change-aware} \\
\bottomrule
\end{tabular}
\end{table}

\begin{table}[H]\centering\small
\caption{知识点总表}
\begin{tabular}{clll}
\toprule
\# & 知识点 & 核心要点 & 难度 \\
\midrule
1 & 放松静态假设 & 假设藏在“重观测约束位姿”与“多帧融合”两处 & $\bigstar\bigstar$ \\
2 & 三轴刻画 & 动态类型 / 理解细节 / 在线离线 & $\bigstar\bigstar$ \\
3 & 短/长期相对率定义 & 变化率/观测率之比；观测者属性 & $\bigstar\bigstar\bigstar$ \\
4 & 动态外点剔除 & 冗余 $+$ 鲁棒核 $+$ 多视一致 $+$ 语义先验 & $\bigstar\bigstar$ \\
5 & “可动”!=“在动” & 语义须叠多视几何检验 & $\bigstar\bigstar$ \\
6 & 动态物体因子图 & 增广物体位姿/速度 $+$ 恒速因子；接母方程 & $\bigstar\bigstar\bigstar\bigstar$ \\
7 & 稠密动态 & 一物一稠密模型（多 TSDF）& $\bigstar\bigstar$ \\
8 & NRSfM 与浮动地图歧义 & 刚/非刚分离病态；尺度歧义的近亲 & $\bigstar\bigstar\bigstar$ \\
9 & ED 图蒙皮 & 稀疏节点变换插值，自由度大降 & $\bigstar\bigstar\bigstar$ \\
10 & ARAP / Killing 能量 & 不拉伸 / 近等距；自补能量泛函 & $\bigstar\bigstar\bigstar\bigstar$ \\
11 & 非刚性融合 & TSDF/surfel $+$ 形变场（复用 ch:dense）& $\bigstar\bigstar\bigstar$ \\
12 & 无证据 vs 缺席证据 & 占据三态方能区分；驱动变化检测 & $\bigstar\bigstar\bigstar$ \\
13 & 边缘化 baking-in & 拓扑冻结；歧义因子须延迟边缘化 & $\bigstar\bigstar\bigstar$ \\
14 & 局部一致性 & change-aware 的可识别性原理 & $\bigstar\bigstar\bigstar\bigstar$ \\
15 & 物体实例重定位 & FPFH/SHOT vs 学习式；$\SE(3)$ 不变识别·等变配准 & $\bigstar\bigstar\bigstar$ \\
\bottomrule
\end{tabular}
\end{table}

% ════════════════════════════════════════════════════════════════
\section*{延伸阅读}
\begin{itemize}
  \item \textbf{本章蓝本}：SLAM Handbook\cite{schmid2026dynamic} 第 15 章《Dynamic and Deformable SLAM》——动态效应刻画、短/长期动态、可变形 SLAM 的系统综述（本章框架与系统选择的出处）。
  \item \textbf{动态物体剔除与跟踪}：DynaSLAM\cite{bescos2018dynaslam}（语义 $+$ 多视几何剔除 $+$ 背景修补）；联合相机运动与场景流\cite{jaimez2017fast}；稠密动态重建 Co-Fusion\cite{runz2017cofusion}。学习式动态分割 MonoRec\cite{wimbauer2021monorec}（\cref{ch:learning} 已引）。
  \item \textbf{可变形稠密融合}：DynamicFusion\cite{newcombe2015dynamicfusion}（canonical map $+$ ED $+$ ARAP）、KillingFusion\cite{slavcheva2017killingfusion}（Killing 正则、拓扑变化）；其复用的 TSDF/surfel 见 \cref{ch:dense}（\cref{sec:dense-tsdf,sec:dense-surfel}）。
  \item \textbf{形变能量（本章自补的数学根）}：嵌入式形变图\cite{sumner2007embedded}、as-rigid-as-possible 表面建模\cite{sorkine2007arap}——几何处理的经典文献，本章 \cref{eq:def-skinning,eq:def-arap,eq:def-killing} 的出处。
  \item \textbf{长期/change-aware 与回环}：长期定位、变化检测、局部一致性的概念框架见\cite{schmid2026dynamic};边缘化的代价与稀疏求解见本书 \cref{ch:nlopt}（\cref{eq:info-marginal}）;地点识别与回环见 \cref{ch:loop}。
\end{itemize}

\section*{本章与前后章节的关系}
\begin{table}[H]\centering\small
\begin{tabular}{lll}
\toprule
章节 & 关系 & 衔接点 \\
\midrule
\cref{ch:intro} 初识 SLAM & 提供静态假设与动态预告 & \cref{sec:intro-dynamic} 是本章的锚；\cref{eq:intro-nls} 母方程 \\
\cref{ch:nlopt} 非线性优化 & 因子图、鲁棒核、边缘化、交替优化 & 动态因子图照此求解；边缘化代价 \cref{eq:info-marginal} \\
\cref{ch:dense} 稠密建图 & TSDF/surfel、占据三态 & 非刚性融合复用之；三态区分“无证据/缺席证据” \\
\cref{ch:loop} 回环检测 & 地点识别、长期数据关联 & 终身定位、物体重定位与之同源 \\
\cref{ch:vo} 视觉里程计 & 重投影、RANSAC、光流 & 改进重投影同构；动态剔除复用鲁棒估计 \\
\cref{ch:learning} 学习赋能 & 网络出掩膜/先验/等变描述子 & 本章是它之后“放松世界假设”的一章 \\
\cref{ch:semantic} 度量-语义 & 语义为“会不会动”供先验 & 语义变化检测、语义一致性接此 \\
\cref{ch:openworld} 开放世界 & 物体可任意增删 & 时序场景理解、世界模型的方向 \\
\bottomrule
\end{tabular}
\end{table}

% ════════════════════════════════════════════════════════════════
\section*{习题}
\begin{enumerate}
  \item \textbf{静态假设藏在哪。}\;指出经典 SLAM 中“静态世界”假设具体藏身的\emph{两处}（提示：回看 \cref{sec:dynamic-frame} 对 \cref{ch:loop} 回环与 \cref{eq:dense-tsdf-fuse} 融合的分析），并各举一个该假设失效后会出现的具体失败现象。
  \item \textbf{短期还是长期。}\;对下列情形判断属短期还是长期动态，并用 \cref{def:dyn-short-long} 的相对率说明理由：(a) 网球被以 340\,Hz 拍摄、球速 250\,km/h；(b) 机器人每天巡逻一次大楼，某区域天天大改；(c) 机器人转身去拿东西、回来发现桌上的瓶子位置变了。
  \item \textbf{动态物体因子图。}\;把 \cref{eq:dyn-reproj} 与 \cref{tab:factors} 的静态观测残差逐项对应：哪个量对应静态地图点、多出来的 $\mathbf{T}_{wi}^{k}$ 起什么作用？再论证“令所有 $\mathbf{T}_{wi}^{k}\equiv\mathbf{I}$ 时 \cref{eq:dyn-map} 退回静态 SLAM”。
  \item \textbf{恒速因子的雅可比。}\;取本书右扰动 $\mathbf{T}\,\mathrm{Exp}(\delta\boldsymbol\xi)$，写出改进重投影残差 \cref{eq:dyn-reproj} 对\emph{相机位姿} $\mathbf{T}_{wi}$ 与\emph{物体位姿} $\mathbf{T}_{wi}^{k}$ 的雅可比各自的链式结构（不必展开到分量，指出与 \cref{ch:vo} 重投影对位姿求导的同构即可）。
  \item \textbf{浮动地图歧义。}\;用一句话解释“浮动地图歧义”，并说明它与 \cref{thm:intro-scale} 单目尺度歧义的相似与不同。要消解它，\cref{sec:dynamic-energy} 的哪些先验起作用？
  \item \textbf{ARAP 为何减旋转。}\;朴素地罚“边长变化”会有什么问题（提示：纯旋转/弯曲）？\cref{eq:def-arap} 为每个节点引入局部旋转 $\mathbf{R}_m$ 后如何修正？写出交替最小化里 $\mathbf{R}_m$ 的闭式解依据（与 \cref{ch:vo} 哪个对齐问题同构）。
  \item \textbf{ARAP vs Killing。}\;比较 \cref{eq:def-arap} 与 \cref{eq:def-killing}：它们各罚什么量？为什么 Killing 正则对\emph{拓扑变化}（如双手合拢）更友好，而 ARAP 不然？
  \item \textbf{非刚性融合复用了什么。}\;对照 \cref{eq:def-fusion-energy} 与 \cref{ch:dense} 的静态 KinectFusion，列出 DynamicFusion \emph{照搬}了 TSDF 的哪几步、\emph{只改}了哪一步。\cref{eq:dense-tsdf-fuse} 的 $W_\eta$ 封顶在这里起什么作用？
  \item \textbf{无证据 vs 缺席证据。}\;用一个具体场景说明二者之别，并解释为什么\emph{占据/TSDF}（三态）能区分、而\emph{稠密点云}（\cref{sec:dense-rgbd-pc}）不能。这对变化检测意味着什么？
  \item \textbf{（选读）边缘化的代价。}\;为什么对\emph{尚有歧义的回环}延迟边缘化很重要？把它与 \cref{eq:info-marginal} 的舒尔补、\cref{ch:nlopt} 的“拓扑不可改”联系起来，论证“边缘化保信息但损灵活性”。
  \item \textbf{（选读）局部一致性。}\;用自己的话陈述 change-aware SLAM 的可识别性原理（局部一致性），并说明它与 \cref{sec:dynamic-removal} 的“多视一致性”、\cref{ch:loop} 的“几何验证回环”共享的是哪一条思路。
\end{enumerate}
